% =====================================================================
%  glossaire.tex : entrées du glossaire (chargé par main.tex via % =====================================================================
%  glossaire.tex : entrées du glossaire (chargé par main.tex via % =====================================================================
%  glossaire.tex : entrées du glossaire (chargé par main.tex via % =====================================================================
%  glossaire.tex : entrées du glossaire (chargé par main.tex via \input{glossaire})
%  Utilisation dans le texte : \gls{iaas} (le mot apparaît en rouge) ; \Gls{...} met la
%  première lettre en capitale.
%  Toutes les entrées sont imprimées en fin de document (\glsaddall dans main.tex).
%  La clé sort={...} donne le texte brut utilisé pour le tri (accents, commandes).
%  Compilation : deux passes de pdflatex suffisent (\makenoidxglossaries).
% =====================================================================

\newglossaryentry{apirest}{name={API REST}, sort={API REST},
  description={Interface d'accès à un service par requêtes HTTP ; chaque service OpenStack expose la sienne, et la ligne de commande comme Horizon n'en sont que des clients}}

\newglossaryentry{appcred}{name={application credential}, sort={application credential},
  description={Identifiant et secret délivrés à une application pour qu'elle s'authentifie à la place de l'utilisateur, avec seulement une partie de ses rôles et, éventuellement, une date d'expiration}}

\newglossaryentry{authentification}{name={authentification}, sort={authentification},
  description={Processus qui confirme l'identité d'un utilisateur ; Keystone valide les identifiants puis délivre un token}}

\newglossaryentry{autorisation}{name={autorisation}, sort={autorisation},
  description={Décision d'autoriser ou de refuser une action à un utilisateur authentifié, d'après ses rôles ; chaque service OpenStack l'applique avec ses propres règles}}

\newglossaryentry{cadf}{name={CADF}, sort={CADF},
  description={\emph{Cloud Auditing Data Federation} : format normalisé d'événements d'audit, utilisé par défaut pour les notifications de Keystone}}

\newglossaryentry{catalogue}{name={catalogue de services}, sort={catalogue de services}, text={catalogue},
  description={Liste des services de la plateforme et de leurs endpoints, par région et par interface ; Keystone l'inclut dans les tokens limités à un projet}}

\newglossaryentry{domaine}{name={domaine}, sort={domaine},
  description={Conteneur de plus haut niveau de Keystone : il regroupe des projets, des utilisateurs et des groupes qui forment une même frontière administrative}}

\newglossaryentry{endpoint}{name={endpoint}, sort={endpoint},
  description={Adresse réseau, généralement une URL, à laquelle on accède à un service ; Keystone en distingue trois interfaces : \texttt{public}, \texttt{internal} et \texttt{admin}}}

\newglossaryentry{federation}{name={fédération d'identité}, sort={federation d'identite},
  description={Délégation de l'authentification à un fournisseur d'identité externe de confiance, par SAML 2.0 ou OpenID Connect}}

\newglossaryentry{fernet}{name={Fernet}, sort={Fernet},
  description={Format de token par défaut de Keystone : chaîne opaque chiffrée (AES-128-CBC) et signée (HMAC-SHA256) avec des clés conservées dans un dépôt local}}

\newglossaryentry{groupe}{name={groupe}, sort={groupe},
  description={Ensemble d'utilisateurs d'un domaine ; un rôle attribué à un groupe s'applique à tous ses membres}}

\newglossaryentry{hyperviseur}{name={hyperviseur}, sort={hyperviseur},
  description={Logiciel qui fait tourner plusieurs machines virtuelles isolées sur un même serveur physique (par exemple KVM) ; OpenStack le pilote sans le remplacer}}

\newglossaryentry{iaas}{name={IaaS}, sort={IaaS},
  description={\emph{Infrastructure as a Service} : modèle de cloud où l'infrastructure (calcul, réseau, stockage) est fournie à la demande ; l'utilisateur gère ce qui s'exécute dessus}}

\newglossaryentry{idp}{name={fournisseur d'identité (IdP)}, sort={fournisseur d'identite}, text={fournisseur d'identité}, plural={fournisseurs d'identité},
  description={Service externe qui authentifie l'utilisateur et transmet à Keystone une assertion signée contenant son identité et ses attributs}}

\newglossaryentry{jws}{name={JWS}, sort={JWS},
  description={\emph{JSON Web Signature} : second format de token de Keystone, signé avec une clé asymétrique (ES256) ; son contenu n'est pas opaque}}

\newglossaryentry{ldap}{name={LDAP}, sort={LDAP},
  description={\emph{Lightweight Directory Access Protocol} : protocole d'accès à un annuaire (par exemple Active Directory) ; Keystone peut y lire utilisateurs et groupes, en lecture seule}}

\newglossaryentry{mapping}{name={mapping}, sort={mapping},
  description={Ensemble de règles qui traduisent les attributs fournis par un fournisseur d'identité en utilisateur, groupes et projets Keystone}}

\newglossaryentry{mfa}{name={authentification multifacteur (MFA)}, sort={authentification multifacteur}, text={authentification multifacteur},
  description={Authentification qui exige plusieurs méthodes successives, par exemple un mot de passe puis un code TOTP}}

\newglossaryentry{middleware}{name={middleware}, sort={middleware},
  description={Composant WSGI placé devant l'application d'un service ; \texttt{keystonemiddleware} (module \texttt{auth\_token}) y valide le token de chaque requête auprès de Keystone}}

\newglossaryentry{oidc}{name={OpenID Connect (OIDC)}, sort={OpenID Connect}, text={OpenID Connect},
  description={Protocole de fédération au format JSON, construit sur OAuth 2.0}}

\newglossaryentry{policy}{name={policy}, sort={policy}, plural={policies},
  description={Règle d'accès évaluée par chaque service (bibliothèque \texttt{oslo.policy}) pour décider si les rôles présents dans le token autorisent l'action demandée}}

\newglossaryentry{portee}{name={portée (scope)}, sort={portee}, text={portée}, plural={portées},
  description={Périmètre sur lequel un token confère des droits : un projet, un domaine ou le système entier ; un token non délimité n'a aucune portée}}

\newglossaryentry{projet}{name={projet (tenant)}, sort={projet}, text={projet}, plural={projets},
  description={Conteneur qui regroupe ou isole des ressources ; « tenant » est l'ancien nom du projet}}

\newglossaryentry{rbac}{name={RBAC}, sort={RBAC},
  description={\emph{Role-Based Access Control} : contrôle d'accès fondé sur les rôles attribués aux utilisateurs et aux groupes}}

\newglossaryentry{region}{name={région}, sort={region},
  description={Division générale d'un déploiement OpenStack, qui peut contenir des sous-régions ; chaque endpoint du catalogue est rattaché à une région}}

\newglossaryentry{role}{name={rôle}, sort={role},
  description={Ensemble de droits attribué à un utilisateur ou à un groupe sur un projet, un domaine ou le système ; les services l'interprètent pour décider des accès}}

\newglossaryentry{saml2}{name={SAML 2.0}, sort={SAML},
  description={\emph{Security Assertion Markup Language} : protocole de fédération au format XML, courant dans les universités et les entreprises}}

\newglossaryentry{sp}{name={fournisseur de services (SP)}, sort={fournisseur de services}, text={fournisseur de services}, plural={fournisseurs de services},
  description={Dans une fédération, service qui s'appuie sur un fournisseur d'identité ; ici Keystone, dont les tokens servent aux autres services OpenStack}}

\newglossaryentry{sso}{name={SSO}, sort={SSO},
  description={\emph{Single Sign-On}, authentification unique : une seule connexion donne accès à plusieurs services}}

\newglossaryentry{token}{name={token (jeton)}, sort={token}, text={token}, plural={tokens},
  description={Chaîne de caractères qui prouve l'identité d'un utilisateur et donne accès aux API ; valide pour une durée limitée, elle peut être révoquée à tout moment}}

\newglossaryentry{totp}{name={TOTP}, sort={TOTP},
  description={\emph{Time-based One-Time Password} : code à usage unique de six chiffres, calculé à partir d'un secret partagé et de l'heure ; second facteur pris en charge par Keystone}}

\newglossaryentry{trust}{name={trust}, sort={trust},
  description={Délégation par un utilisateur (le \emph{trustor}) d'une partie de ses rôles sur un projet à un autre utilisateur (le \emph{trustee})}}

\newglossaryentry{utilisateur}{name={utilisateur}, sort={utilisateur},
  description={Représentation numérique d'une personne, d'un système ou d'un service qui utilise les services OpenStack}}

\newglossaryentry{wsgi}{name={WSGI}, sort={WSGI},
  description={\emph{Web Server Gateway Interface} : interface standard entre un serveur web et une application Python ; Keystone s'exécute derrière Apache grâce à \texttt{mod\_wsgi}}}

% ---- Nova et Placement ------------------------------------------------
\newglossaryentry{instance}{name={instance}, sort={instance},
  description={Machine virtuelle créée et gérée par Nova ; on parle aussi de serveur (\texttt{openstack server})}}

\newglossaryentry{flavor}{name={gabarit (flavor)}, sort={gabarit}, text={gabarit},
  description={Modèle de taille d'une instance : nombre de vCPU, mémoire, disque racine, disque éphémère, swap, et propriétés supplémentaires (\emph{extra specs})}}

\newglossaryentry{noeudcalcul}{name={nœud de calcul (compute node)}, sort={noeud de calcul}, text={nœud de calcul},
  description={Serveur physique qui héberge les instances : il exécute \texttt{nova-compute} et un hyperviseur}}

\newglossaryentry{ordonnanceur}{name={ordonnanceur (scheduler)}, sort={ordonnanceur}, text={ordonnanceur},
  description={Composant de Nova (\texttt{nova-scheduler}) qui choisit le nœud de calcul où placer une instance}}

\newglossaryentry{conductor}{name={conductor}, sort={conductor},
  description={Composant de Nova (\texttt{nova-conductor}) qui coordonne les opérations longues (construction, redimensionnement, migration) et sert d'intermédiaire entre les nœuds de calcul et la base de données}}

\newglossaryentry{cellule}{name={cellule (cell)}, sort={cellule}, text={cellule},
  description={Groupe de nœuds de calcul de Nova qui possède sa propre base de données et sa propre file de messages ; tout déploiement a au moins une cellule}}

\newglossaryentry{libvirt}{name={libvirt}, sort={libvirt},
  description={Bibliothèque et service de pilotage des hyperviseurs (KVM/QEMU en particulier) que \texttt{nova-compute} utilise pour créer et contrôler les machines virtuelles}}

\newglossaryentry{busmessages}{name={bus de messages}, sort={bus de messages}, text={bus de messages},
  description={File de messages (RabbitMQ en général) par laquelle les composants d'un service échangent des appels distants (RPC, via \texttt{oslo.messaging})}}

\newglossaryentry{microversion}{name={microversion}, sort={microversion},
  description={Numéro de version fin d'une API OpenStack (pour Nova : 2.1, 2.2, 2.3 et ainsi de suite) : le client demande la version dont il a besoin, ce qui garde l'API compatible}}

\newglossaryentry{migrationchaud}{name={migration à chaud (live migration)}, sort={migration a chaud}, text={migration à chaud},
  description={Déplacement d'une instance en cours d'exécution d'un nœud de calcul vers un autre, sans l'arrêter}}

\newglossaryentry{surallocation}{name={surallocation (overcommit)}, sort={surallocation}, text={surallocation},
  description={Attribution de plus de ressources virtuelles qu'il n'y en a de physiques ; elle se règle par un ratio d'allocation}}

\newglossaryentry{agregat}{name={agrégat d'hôtes (host aggregate)}, sort={agregat}, text={agrégat d'hôtes},
  description={Groupe de nœuds de calcul défini par l'administrateur, avec des métadonnées ; il sert à l'ordonnancement et à définir des zones de disponibilité}}

\newglossaryentry{zonedispo}{name={zone de disponibilité}, sort={zone de disponibilite}, text={zone de disponibilité},
  description={Groupe de nœuds de calcul que l'utilisateur peut viser à la création d'une instance (par exemple une salle ou une baie) ; c'est un agrégat d'hôtes visible des utilisateurs}}

\newglossaryentry{fournisseur}{name={fournisseur de ressources (resource provider)}, sort={fournisseur de ressources}, text={fournisseur de ressources},
  description={Entité de Placement qui offre un inventaire de ressources consommables, par exemple un nœud de calcul ou un pool de stockage partagé}}

\newglossaryentry{inventaire}{name={inventaire}, sort={inventaire},
  description={Quantité d'une classe de ressources offerte par un fournisseur : total, réservé, ratio d'allocation, taille minimale, maximale et pas d'allocation}}

\newglossaryentry{classeressources}{name={classe de ressources}, sort={classe de ressources}, text={classe de ressources},
  description={Type de ressource quantitative suivie par Placement : \texttt{VCPU}, \texttt{MEMORY\_MB}, \texttt{DISK\_GB}, ou une classe personnalisée préfixée par \texttt{CUSTOM\_}}}

\newglossaryentry{allocation}{name={allocation}, sort={allocation},
  description={Quantité de ressources d'un fournisseur attribuée à un consommateur ; elle diminue d'autant la capacité disponible}}

\newglossaryentry{consommateur}{name={consommateur (consumer)}, sort={consommateur}, text={consommateur},
  description={Entité qui détient des allocations dans Placement, identifiée par un UUID : pour Nova, une instance ou une migration}}

\newglossaryentry{trait}{name={trait}, sort={trait},
  description={Caractéristique qualitative d'un fournisseur (par exemple \texttt{HW\_CPU\_X86\_AVX2} ou \texttt{STORAGE\_DISK\_SSD}) ; les traits non standard commencent par \texttt{CUSTOM\_}}}

\newglossaryentry{claim}{name={réservation (claim)}, sort={reservation}, text={réservation},
  description={Création d'une allocation dans Placement pour une instance, faite par l'ordonnanceur une fois l'hôte choisi ; si elle échoue, c'est qu'un autre consommateur a pris les ressources entre-temps}}

% ---- Glance ----------------------------------------------------------
\newglossaryentry{image}{name={image}, sort={image},
  description={Fichier disque d'une machine virtuelle (système installé) qu'on duplique pour créer des instances ; Glance en conserve les données et les métadonnées}}

\newglossaryentry{qcow2}{name={qcow2}, sort={qcow2},
  description={Format de disque de QEMU à copie sur écriture : le fichier ne grossit qu'avec les données écrites et peut s'appuyer sur un fichier de base}}

\newglossaryentry{store}{name={store (backend de stockage)}, sort={store}, text={store},
  description={Système où Glance garde les données des images : système de fichiers, Ceph (rbd), Swift, Cinder, VMware ; en lecture seule pour http}}

% ---- Neutron ---------------------------------------------------------
\newglossaryentry{reseau}{name={réseau (network)}, sort={reseau}, text={réseau}, plural={réseaux},
  description={Dans Neutron, domaine de diffusion de niveau 2 isolé, auquel on rattache des ports}}

\newglossaryentry{sousreseau}{name={sous-réseau (subnet)}, sort={sous-reseau}, text={sous-réseau}, plural={sous-réseaux},
  description={Plage d'adresses IP associée à un réseau Neutron, avec sa passerelle et son serveur DHCP éventuel}}

\newglossaryentry{port}{name={port}, sort={port},
  description={Point de connexion d'un équipement (en général la carte réseau d'une instance) à un réseau Neutron ; il porte une adresse MAC, des adresses IP et des groupes de sécurité}}

\newglossaryentry{routeur}{name={routeur (router)}, sort={routeur}, text={routeur},
  description={Objet Neutron qui relie des réseaux entre eux et vers l'extérieur, avec traduction d'adresses (NAT)}}

\newglossaryentry{groupesecurite}{name={groupe de sécurité}, sort={groupe de securite}, text={groupe de sécurité}, plural={groupes de sécurité},
  description={Ensemble de règles de pare-feu appliqué aux ports : par défaut tout le trafic entrant est refusé et tout le trafic sortant autorisé ; les règles sont avec état}}

\newglossaryentry{ipflottante}{name={IP flottante (floating IP)}, sort={IP flottante}, text={IP flottante},
  description={Adresse publique associée à la demande à un port d'instance, et traduite en son adresse privée par le routeur}}

\newglossaryentry{provider}{name={réseau provider}, sort={reseau provider}, text={réseau provider},
  description={Réseau créé par l'administrateur et rattaché directement à un réseau physique existant (flat ou VLAN)}}

\newglossaryentry{selfservice}{name={réseau self-service}, sort={reseau self-service}, text={réseau self-service},
  description={Réseau virtuel (overlay) créé par un projet sans droits d'administrateur ; il a besoin d'un routeur pour sortir}}

\newglossaryentry{overlay}{name={overlay (réseau superposé)}, sort={overlay}, text={overlay},
  description={Réseau virtuel dont les trames sont encapsulées (VXLAN, Geneve) dans des paquets IP du réseau physique}}

\newglossaryentry{vlan}{name={VLAN}, sort={VLAN},
  description={Réseau local virtuel 802.1Q identifié par un numéro de 1 à 4094 ; la limite de 4094 réseaux restreint son usage pour les projets}}

\newglossaryentry{vxlan}{name={VXLAN, Geneve}, sort={VXLAN}, text={VXLAN},
  description={Protocoles de tunnel qui encapsulent des trames Ethernet dans des paquets UDP/IP ; ils servent aux réseaux overlay et n'ont pas la limite des VLAN}}

\newglossaryentry{ml2}{name={ML2 (Modular Layer 2)}, sort={ML2}, text={ML2},
  description={Plugin de base de Neutron : il orchestre des \emph{type drivers} (flat, VLAN, VXLAN\dots) et des \emph{mechanism drivers} (Open vSwitch, OVN\dots)}}

\newglossaryentry{ovs}{name={Open vSwitch (OVS)}, sort={Open vSwitch}, text={Open vSwitch},
  description={Commutateur virtuel programmable : il relie sur chaque nœud les interfaces des instances, les tunnels et les interfaces physiques, selon des règles (\emph{flows}) OpenFlow}}

\newglossaryentry{ovn}{name={OVN (Open Virtual Network)}, sort={OVN}, text={OVN},
  description={Système de réseau virtuel bâti sur Open vSwitch, avec deux bases de données et un démon par nœud ; il distribue le routage, le DHCP et les règles de sécurité}}

\newglossaryentry{agentneutron}{name={agent Neutron}, sort={agent Neutron}, text={agent},
  description={Processus qui, sur un nœud, traduit les objets Neutron en configuration réelle : agent L2 (Open vSwitch), L3 (routeurs), DHCP, métadonnées}}

\newglossaryentry{nat}{name={NAT}, sort={NAT},
  description={Traduction d'adresses réseau : le routeur remplace l'adresse source ou destination d'un paquet (SNAT pour la sortie, DNAT pour l'entrée)}}

% ---- Cinder -----------------------------------------------------------
\newglossaryentry{volume}{name={volume}, sort={volume},
  description={Disque virtuel persistant fourni par Cinder : il a une taille et un type, et peut être attaché à une instance, détaché, puis rattaché ailleurs sans perdre ses données}}

\newglossaryentry{snapshot}{name={snapshot}, sort={snapshot},
  description={Image ponctuelle d'un volume, conservée sur le même backend ; elle sert au retour arrière et au clonage, mais disparaît avec le backend}}

\newglossaryentry{sauvegarde}{name={sauvegarde (backup)}, sort={sauvegarde}, text={sauvegarde},
  description={Copie complète ou incrémentale d'un volume vers un dépôt distinct (Swift par défaut) ; elle survit à la perte du backend}}

\newglossaryentry{typevolume}{name={type de volume}, sort={type de volume}, plural={types de volume},
  description={Étiquette qui porte des propriétés (\emph{extra specs}) : backend visé, qualité de service, accès multiple ; Cinder s'en sert pour choisir le backend d'un volume}}

\newglossaryentry{backendcinder}{name={backend (Cinder)}, sort={backend}, text={backend}, plural={backends},
  description={Système de stockage qui héberge réellement les volumes (LVM, Ceph, NFS, baie SAN) ; \texttt{cinder-volume} le pilote par un driver}}

\newglossaryentry{qos}{name={QoS (qualité de service)}, sort={QoS}, text={QoS},
  description={Limite de débit appliquée à un volume, en opérations ou en octets par seconde ; elle s'associe à un type de volume}}

\newglossaryentry{osbrick}{name={os-brick}, sort={os-brick},
  description={Bibliothèque qui connecte un hôte à un volume (iSCSI, RBD, NFS\dots) ; Nova l'utilise sur le nœud de calcul pour attacher les volumes}}

% ---- Horizon ----------------------------------------------------------
\newglossaryentry{django}{name={Django}, sort={Django},
  description={Cadriciel Web en Python sur lequel est écrit Horizon ; il fournit les sessions, l'authentification par des backends et le système de gabarits}}

\newglossaryentry{dashboard}{name={dashboard (tableau de bord)}, sort={dashboard}, text={dashboard},
  description={Entrée de la navigation principale d'Horizon (\emph{Project}, \emph{Admin}, \emph{Identity}, \emph{Settings}) ; il regroupe des groupes de panneaux}}

\newglossaryentry{panneau}{name={panneau (panel)}, sort={panneau}, text={panneau}, plural={panneaux},
  description={Page de la navigation latérale d'un dashboard, déclarée dans un fichier \texttt{panel.py} ; elle a ses routes, ses vues et ses gabarits}}

\newglossaryentry{memcached}{name={memcached}, sort={memcached},
  description={Service de cache en mémoire, partagé entre processus et serveurs ; Horizon l'utilise couramment pour conserver les sessions}}

% ---- Octavia ----------------------------------------------------------
\newglossaryentry{loadbalancer}{name={répartiteur de charge (load balancer)}, sort={repartiteur de charge}, text={répartiteur de charge}, plural={répartiteurs de charge},
  description={Service qui reçoit le trafic sur une adresse unique (la VIP) et le distribue entre plusieurs serveurs ; c'est l'objet principal que l'utilisateur crée dans Octavia}}

\newglossaryentry{listener}{name={listener}, sort={listener}, plural={listeners},
  description={Point d'écoute d'un répartiteur : un protocole et un port (par exemple HTTP, 80) sur la VIP ; il envoie le trafic reçu vers un pool}}

\newglossaryentry{poolocta}{name={pool (Octavia)}, sort={pool}, text={pool}, plural={pools},
  description={Groupe de serveurs qui reçoivent le trafic d'un listener, avec un algorithme de répartition et, au besoin, un health monitor}}

\newglossaryentry{memberocta}{name={member (Octavia)}, sort={member}, text={member}, plural={members},
  description={Serveur d'un pool, repéré par une adresse IP et un port ; il peut être une instance Nova ou n'importe quelle adresse joignable}}

\newglossaryentry{healthmonitor}{name={health monitor}, sort={health monitor},
  description={Sonde que le répartiteur envoie à chaque member (TCP, HTTP, HTTPS, PING\dots) ; un member qui échoue est retiré de la distribution jusqu'à ce qu'il réponde de nouveau}}

\newglossaryentry{amphore}{name={amphore}, sort={amphore}, plural={amphores},
  description={Instance Nova, créée par Octavia, qui exécute HAProxy et porte la VIP ; c'est elle qui écoute et répartit réellement le trafic}}

\newglossaryentry{vip}{name={VIP (adresse virtuelle)}, sort={VIP}, text={VIP},
  description={Adresse IP du répartiteur de charge, celle que les clients utilisent ; elle est portée par un port Neutron et peut recevoir une adresse flottante}}

\newglossaryentry{l7}{name={politique L7 (L7 policy)}, sort={politique L7}, text={politique L7}, plural={politiques L7},
  description={Règle de routage au niveau applicatif (chemin d'URL, en-tête, nom d'hôte\dots) qui redirige une requête vers un autre pool ou vers une URL}}

\newglossaryentry{haproxy}{name={HAProxy}, sort={HAProxy},
  description={Logiciel libre de répartition de charge TCP et HTTP ; l'amphore de référence d'Octavia l'exécute}}

% ---- Swift ------------------------------------------------------------
\newglossaryentry{objetswift}{name={objet (Swift)}, sort={objet}, text={objet}, plural={objets},
  description={Donnée stockée dans Swift : le contenu d'un fichier et ses métadonnées, adressé par son chemin \texttt{/compte/conteneur/objet}}}

\newglossaryentry{conteneur}{name={conteneur (Swift)}, sort={conteneur}, text={conteneur}, plural={conteneurs},
  description={Espace de noms d'un compte Swift qui regroupe des objets ; il porte les droits d'accès, la politique de stockage et d'autres métadonnées}}

\newglossaryentry{proxyswift}{name={proxy (Swift)}, sort={proxy}, text={proxy}, plural={proxys},
  description={Service \texttt{swift-proxy-server} : seul point d'entrée de l'API, il authentifie, consulte le ring et relaie les requêtes vers les serveurs de stockage}}

\newglossaryentry{ring}{name={ring}, sort={ring}, plural={rings},
  description={Fichier qui associe chaque partition à des disques précis ; Swift en tient un pour les comptes, un pour les conteneurs et un par politique de stockage d'objets}}

\newglossaryentry{partition}{name={partition (Swift)}, sort={partition}, text={partition}, plural={partitions},
  description={Unité de répartition du ring : un intervalle des empreintes MD5 des chemins ; le ring place ses répliques sur des disques de zones différentes}}

\newglossaryentry{handoff}{name={handoff}, sort={handoff}, plural={handoffs},
  description={Disque de remplacement choisi par le ring lorsqu'un disque normal est en panne ; la réplication ramène ensuite les données à leur place}}

\newglossaryentry{erasurecoding}{name={codage d'effacement (erasure coding)}, sort={codage d'effacement}, text={codage d'effacement},
  description={Méthode de redondance qui découpe un objet en fragments de données et de parité ; Swift l'offre par une politique de stockage, pour un surcoût moindre que la triple réplication}}

\newglossaryentry{politiquestockage}{name={politique de stockage (storage policy)}, sort={politique de stockage}, text={politique de stockage}, plural={politiques de stockage},
  description={Choix, par conteneur, d'un ring d'objets et donc du nombre de répliques, du type de disques ou du codage d'effacement}}

% ---- Heat -------------------------------------------------------------
\newglossaryentry{modeleheat}{name={modèle (template)}, sort={modele}, text={modèle}, plural={modèles},
  description={Fichier qui décrit l'infrastructure voulue (ressources, paramètres, sorties) ; Heat le lit pour créer une stack. Son format natif est HOT, en YAML}}

\newglossaryentry{stack}{name={stack}, sort={stack}, plural={stacks},
  description={Ensemble de ressources créé par Heat à partir d'un modèle ; c'est l'unité qu'on crée, met à jour et supprime}}

\newglossaryentry{ressourceheat}{name={ressource (Heat)}, sort={ressource}, text={ressource}, plural={ressources},
  description={Élément d'infrastructure d'une stack (réseau, serveur, volume\dots), instancié d'après un type de la forme \texttt{OS::Service::Type}}}

\newglossaryentry{hot}{name={HOT (Heat Orchestration Template)}, sort={HOT}, text={HOT},
  description={Format de modèle natif de Heat, écrit en YAML ; sa version (\texttt{heat\_template\_version}) détermine les fonctions disponibles}}

\newglossaryentry{environnementheat}{name={environnement (Heat)}, sort={environnement}, text={environnement},
  description={Fichier passé avec \texttt{-e} qui fournit des valeurs de paramètres et peut remplacer des types de ressources (\texttt{resource\_registry})}}

\newglossaryentry{graphedeps}{name={graphe de dépendances}, sort={graphe de dependances}, text={graphe de dépendances},
  description={Graphe orienté que Heat construit à partir des références entre ressources ; il fixe l'ordre de création, et la suppression suit l'ordre inverse}}

% ---- Ironic -----------------------------------------------------------
\newglossaryentry{baremetal}{name={bare metal}, sort={bare metal}, text={bare metal},
  description={Serveur physique utilisé directement par un seul locataire, sans hyperviseur ; Ironic le provisionne et Nova le présente comme une instance}}

\newglossaryentry{noeudironic}{name={nœud (Ironic)}, sort={noeud (Ironic)}, text={nœud}, plural={nœuds},
  description={Enregistrement d'Ironic qui représente un serveur physique : son type de matériel, ses identifiants de gestion, ses ports, sa classe de ressources et son état de provisionnement}}

\newglossaryentry{portironic}{name={port (Ironic)}, sort={port (Ironic)}, text={port}, plural={ports},
  description={Carte réseau d'un nœud, déclarée dans Ironic par son adresse MAC ; à ne pas confondre avec le port Neutron, qui lui est lié lors du raccordement au réseau}}

\newglossaryentry{conductorironic}{name={conductor (Ironic)}, sort={conductor (Ironic)}, text={conductor}, plural={conductors},
  description={Composant d'Ironic (\texttt{ironic-conductor}) qui agit sur les serveurs physiques : alimentation, déploiement, nettoyage, inspection ; chaque conductor gère une partie des nœuds}}

\newglossaryentry{bmc}{name={BMC (contrôleur de gestion)}, sort={BMC}, text={BMC},
  description={\emph{Baseboard Management Controller} : contrôleur intégré au serveur, doté de sa propre adresse réseau ; il reçoit les ordres d'alimentation et de démarrage (IPMI, Redfish)}}

\newglossaryentry{ipmi}{name={IPMI}, sort={IPMI}, text={IPMI},
  description={\emph{Intelligent Platform Management Interface} : protocole ancien de gestion à distance par le BMC (alimentation, périphérique de démarrage) ; Ironic l'utilise avec la commande \texttt{ipmitool}}}

\newglossaryentry{redfish}{name={Redfish}, sort={Redfish}, text={Redfish},
  description={Protocole de gestion à distance des serveurs par API REST, proposé par le BMC ; Ironic l'utilise avec le type de matériel \texttt{redfish}}}

\newglossaryentry{pxe}{name={PXE (démarrage par le réseau)}, sort={PXE}, text={PXE},
  description={\emph{Preboot Execution Environment} : démarrage par le réseau, où un serveur DHCP indique au serveur le programme à charger et un serveur TFTP (ou HTTP) le fournit ; iPXE en est une variante}}

\newglossaryentry{ipa}{name={IPA (Ironic Python Agent)}, sort={IPA}, text={IPA},
  description={Agent Python qui s'exécute dans un ramdisk sur le serveur : il donne au conductor un accès à distance et réalise de l'intérieur l'écriture de l'image, le nettoyage et l'inspection}}

\newglossaryentry{hardwaretype}{name={type de matériel (hardware type)}, sort={type de materiel}, text={type de matériel}, plural={types de matériel},
  description={Ensemble d'interfaces matérielles qui définit comment Ironic pilote une famille de serveurs (\texttt{ipmi}, \texttt{redfish} ou un type propre à un constructeur) ; on parle aussi de pilote (\emph{driver})}}

\newglossaryentry{interfacemat}{name={interface matérielle (hardware interface)}, sort={interface materielle}, text={interface matérielle}, plural={interfaces matérielles},
  description={Brique d'un type de matériel qui gère un aspect du serveur : alimentation, démarrage, déploiement, réseau, inspection, RAID\dots}}

\newglossaryentry{nettoyage}{name={nettoyage (cleaning)}, sort={nettoyage}, text={nettoyage},
  description={Étapes qu'Ironic exécute sur un nœud entre deux locataires, ou avant le premier : effacement des disques, remise du matériel dans une configuration connue ; l'état \texttt{cleaning} les regroupe}}

\newglossaryentry{ramdisk}{name={ramdisk de déploiement}, sort={ramdisk}, text={ramdisk},
  description={Petit système Linux chargé en mémoire par démarrage réseau, qui exécute l'agent IPA pour déployer, nettoyer ou inspecter un nœud}}

\newglossaryentry{etatprov}{name={état de provisionnement (provision state)}, sort={etat de provisionnement}, text={état de provisionnement}, plural={états de provisionnement},
  description={Étape du cycle de vie d'un nœud dans la machine à états d'Ironic (\texttt{enroll}, \texttt{manageable}, \texttt{available}, \texttt{active}\dots)}}
)
%  Utilisation dans le texte : \gls{iaas} (le mot apparaît en rouge) ; \Gls{...} met la
%  première lettre en capitale.
%  Toutes les entrées sont imprimées en fin de document (\glsaddall dans main.tex).
%  La clé sort={...} donne le texte brut utilisé pour le tri (accents, commandes).
%  Compilation : deux passes de pdflatex suffisent (\makenoidxglossaries).
% =====================================================================

\newglossaryentry{apirest}{name={API REST}, sort={API REST},
  description={Interface d'accès à un service par requêtes HTTP ; chaque service OpenStack expose la sienne, et la ligne de commande comme Horizon n'en sont que des clients}}

\newglossaryentry{appcred}{name={application credential}, sort={application credential},
  description={Identifiant et secret délivrés à une application pour qu'elle s'authentifie à la place de l'utilisateur, avec seulement une partie de ses rôles et, éventuellement, une date d'expiration}}

\newglossaryentry{authentification}{name={authentification}, sort={authentification},
  description={Processus qui confirme l'identité d'un utilisateur ; Keystone valide les identifiants puis délivre un token}}

\newglossaryentry{autorisation}{name={autorisation}, sort={autorisation},
  description={Décision d'autoriser ou de refuser une action à un utilisateur authentifié, d'après ses rôles ; chaque service OpenStack l'applique avec ses propres règles}}

\newglossaryentry{cadf}{name={CADF}, sort={CADF},
  description={\emph{Cloud Auditing Data Federation} : format normalisé d'événements d'audit, utilisé par défaut pour les notifications de Keystone}}

\newglossaryentry{catalogue}{name={catalogue de services}, sort={catalogue de services}, text={catalogue},
  description={Liste des services de la plateforme et de leurs endpoints, par région et par interface ; Keystone l'inclut dans les tokens limités à un projet}}

\newglossaryentry{domaine}{name={domaine}, sort={domaine},
  description={Conteneur de plus haut niveau de Keystone : il regroupe des projets, des utilisateurs et des groupes qui forment une même frontière administrative}}

\newglossaryentry{endpoint}{name={endpoint}, sort={endpoint},
  description={Adresse réseau, généralement une URL, à laquelle on accède à un service ; Keystone en distingue trois interfaces : \texttt{public}, \texttt{internal} et \texttt{admin}}}

\newglossaryentry{federation}{name={fédération d'identité}, sort={federation d'identite},
  description={Délégation de l'authentification à un fournisseur d'identité externe de confiance, par SAML 2.0 ou OpenID Connect}}

\newglossaryentry{fernet}{name={Fernet}, sort={Fernet},
  description={Format de token par défaut de Keystone : chaîne opaque chiffrée (AES-128-CBC) et signée (HMAC-SHA256) avec des clés conservées dans un dépôt local}}

\newglossaryentry{groupe}{name={groupe}, sort={groupe},
  description={Ensemble d'utilisateurs d'un domaine ; un rôle attribué à un groupe s'applique à tous ses membres}}

\newglossaryentry{hyperviseur}{name={hyperviseur}, sort={hyperviseur},
  description={Logiciel qui fait tourner plusieurs machines virtuelles isolées sur un même serveur physique (par exemple KVM) ; OpenStack le pilote sans le remplacer}}

\newglossaryentry{iaas}{name={IaaS}, sort={IaaS},
  description={\emph{Infrastructure as a Service} : modèle de cloud où l'infrastructure (calcul, réseau, stockage) est fournie à la demande ; l'utilisateur gère ce qui s'exécute dessus}}

\newglossaryentry{idp}{name={fournisseur d'identité (IdP)}, sort={fournisseur d'identite}, text={fournisseur d'identité}, plural={fournisseurs d'identité},
  description={Service externe qui authentifie l'utilisateur et transmet à Keystone une assertion signée contenant son identité et ses attributs}}

\newglossaryentry{jws}{name={JWS}, sort={JWS},
  description={\emph{JSON Web Signature} : second format de token de Keystone, signé avec une clé asymétrique (ES256) ; son contenu n'est pas opaque}}

\newglossaryentry{ldap}{name={LDAP}, sort={LDAP},
  description={\emph{Lightweight Directory Access Protocol} : protocole d'accès à un annuaire (par exemple Active Directory) ; Keystone peut y lire utilisateurs et groupes, en lecture seule}}

\newglossaryentry{mapping}{name={mapping}, sort={mapping},
  description={Ensemble de règles qui traduisent les attributs fournis par un fournisseur d'identité en utilisateur, groupes et projets Keystone}}

\newglossaryentry{mfa}{name={authentification multifacteur (MFA)}, sort={authentification multifacteur}, text={authentification multifacteur},
  description={Authentification qui exige plusieurs méthodes successives, par exemple un mot de passe puis un code TOTP}}

\newglossaryentry{middleware}{name={middleware}, sort={middleware},
  description={Composant WSGI placé devant l'application d'un service ; \texttt{keystonemiddleware} (module \texttt{auth\_token}) y valide le token de chaque requête auprès de Keystone}}

\newglossaryentry{oidc}{name={OpenID Connect (OIDC)}, sort={OpenID Connect}, text={OpenID Connect},
  description={Protocole de fédération au format JSON, construit sur OAuth 2.0}}

\newglossaryentry{policy}{name={policy}, sort={policy}, plural={policies},
  description={Règle d'accès évaluée par chaque service (bibliothèque \texttt{oslo.policy}) pour décider si les rôles présents dans le token autorisent l'action demandée}}

\newglossaryentry{portee}{name={portée (scope)}, sort={portee}, text={portée}, plural={portées},
  description={Périmètre sur lequel un token confère des droits : un projet, un domaine ou le système entier ; un token non délimité n'a aucune portée}}

\newglossaryentry{projet}{name={projet (tenant)}, sort={projet}, text={projet}, plural={projets},
  description={Conteneur qui regroupe ou isole des ressources ; « tenant » est l'ancien nom du projet}}

\newglossaryentry{rbac}{name={RBAC}, sort={RBAC},
  description={\emph{Role-Based Access Control} : contrôle d'accès fondé sur les rôles attribués aux utilisateurs et aux groupes}}

\newglossaryentry{region}{name={région}, sort={region},
  description={Division générale d'un déploiement OpenStack, qui peut contenir des sous-régions ; chaque endpoint du catalogue est rattaché à une région}}

\newglossaryentry{role}{name={rôle}, sort={role},
  description={Ensemble de droits attribué à un utilisateur ou à un groupe sur un projet, un domaine ou le système ; les services l'interprètent pour décider des accès}}

\newglossaryentry{saml2}{name={SAML 2.0}, sort={SAML},
  description={\emph{Security Assertion Markup Language} : protocole de fédération au format XML, courant dans les universités et les entreprises}}

\newglossaryentry{sp}{name={fournisseur de services (SP)}, sort={fournisseur de services}, text={fournisseur de services}, plural={fournisseurs de services},
  description={Dans une fédération, service qui s'appuie sur un fournisseur d'identité ; ici Keystone, dont les tokens servent aux autres services OpenStack}}

\newglossaryentry{sso}{name={SSO}, sort={SSO},
  description={\emph{Single Sign-On}, authentification unique : une seule connexion donne accès à plusieurs services}}

\newglossaryentry{token}{name={token (jeton)}, sort={token}, text={token}, plural={tokens},
  description={Chaîne de caractères qui prouve l'identité d'un utilisateur et donne accès aux API ; valide pour une durée limitée, elle peut être révoquée à tout moment}}

\newglossaryentry{totp}{name={TOTP}, sort={TOTP},
  description={\emph{Time-based One-Time Password} : code à usage unique de six chiffres, calculé à partir d'un secret partagé et de l'heure ; second facteur pris en charge par Keystone}}

\newglossaryentry{trust}{name={trust}, sort={trust},
  description={Délégation par un utilisateur (le \emph{trustor}) d'une partie de ses rôles sur un projet à un autre utilisateur (le \emph{trustee})}}

\newglossaryentry{utilisateur}{name={utilisateur}, sort={utilisateur},
  description={Représentation numérique d'une personne, d'un système ou d'un service qui utilise les services OpenStack}}

\newglossaryentry{wsgi}{name={WSGI}, sort={WSGI},
  description={\emph{Web Server Gateway Interface} : interface standard entre un serveur web et une application Python ; Keystone s'exécute derrière Apache grâce à \texttt{mod\_wsgi}}}

% ---- Nova et Placement ------------------------------------------------
\newglossaryentry{instance}{name={instance}, sort={instance},
  description={Machine virtuelle créée et gérée par Nova ; on parle aussi de serveur (\texttt{openstack server})}}

\newglossaryentry{flavor}{name={gabarit (flavor)}, sort={gabarit}, text={gabarit},
  description={Modèle de taille d'une instance : nombre de vCPU, mémoire, disque racine, disque éphémère, swap, et propriétés supplémentaires (\emph{extra specs})}}

\newglossaryentry{noeudcalcul}{name={nœud de calcul (compute node)}, sort={noeud de calcul}, text={nœud de calcul},
  description={Serveur physique qui héberge les instances : il exécute \texttt{nova-compute} et un hyperviseur}}

\newglossaryentry{ordonnanceur}{name={ordonnanceur (scheduler)}, sort={ordonnanceur}, text={ordonnanceur},
  description={Composant de Nova (\texttt{nova-scheduler}) qui choisit le nœud de calcul où placer une instance}}

\newglossaryentry{conductor}{name={conductor}, sort={conductor},
  description={Composant de Nova (\texttt{nova-conductor}) qui coordonne les opérations longues (construction, redimensionnement, migration) et sert d'intermédiaire entre les nœuds de calcul et la base de données}}

\newglossaryentry{cellule}{name={cellule (cell)}, sort={cellule}, text={cellule},
  description={Groupe de nœuds de calcul de Nova qui possède sa propre base de données et sa propre file de messages ; tout déploiement a au moins une cellule}}

\newglossaryentry{libvirt}{name={libvirt}, sort={libvirt},
  description={Bibliothèque et service de pilotage des hyperviseurs (KVM/QEMU en particulier) que \texttt{nova-compute} utilise pour créer et contrôler les machines virtuelles}}

\newglossaryentry{busmessages}{name={bus de messages}, sort={bus de messages}, text={bus de messages},
  description={File de messages (RabbitMQ en général) par laquelle les composants d'un service échangent des appels distants (RPC, via \texttt{oslo.messaging})}}

\newglossaryentry{microversion}{name={microversion}, sort={microversion},
  description={Numéro de version fin d'une API OpenStack (pour Nova : 2.1, 2.2, 2.3 et ainsi de suite) : le client demande la version dont il a besoin, ce qui garde l'API compatible}}

\newglossaryentry{migrationchaud}{name={migration à chaud (live migration)}, sort={migration a chaud}, text={migration à chaud},
  description={Déplacement d'une instance en cours d'exécution d'un nœud de calcul vers un autre, sans l'arrêter}}

\newglossaryentry{surallocation}{name={surallocation (overcommit)}, sort={surallocation}, text={surallocation},
  description={Attribution de plus de ressources virtuelles qu'il n'y en a de physiques ; elle se règle par un ratio d'allocation}}

\newglossaryentry{agregat}{name={agrégat d'hôtes (host aggregate)}, sort={agregat}, text={agrégat d'hôtes},
  description={Groupe de nœuds de calcul défini par l'administrateur, avec des métadonnées ; il sert à l'ordonnancement et à définir des zones de disponibilité}}

\newglossaryentry{zonedispo}{name={zone de disponibilité}, sort={zone de disponibilite}, text={zone de disponibilité},
  description={Groupe de nœuds de calcul que l'utilisateur peut viser à la création d'une instance (par exemple une salle ou une baie) ; c'est un agrégat d'hôtes visible des utilisateurs}}

\newglossaryentry{fournisseur}{name={fournisseur de ressources (resource provider)}, sort={fournisseur de ressources}, text={fournisseur de ressources},
  description={Entité de Placement qui offre un inventaire de ressources consommables, par exemple un nœud de calcul ou un pool de stockage partagé}}

\newglossaryentry{inventaire}{name={inventaire}, sort={inventaire},
  description={Quantité d'une classe de ressources offerte par un fournisseur : total, réservé, ratio d'allocation, taille minimale, maximale et pas d'allocation}}

\newglossaryentry{classeressources}{name={classe de ressources}, sort={classe de ressources}, text={classe de ressources},
  description={Type de ressource quantitative suivie par Placement : \texttt{VCPU}, \texttt{MEMORY\_MB}, \texttt{DISK\_GB}, ou une classe personnalisée préfixée par \texttt{CUSTOM\_}}}

\newglossaryentry{allocation}{name={allocation}, sort={allocation},
  description={Quantité de ressources d'un fournisseur attribuée à un consommateur ; elle diminue d'autant la capacité disponible}}

\newglossaryentry{consommateur}{name={consommateur (consumer)}, sort={consommateur}, text={consommateur},
  description={Entité qui détient des allocations dans Placement, identifiée par un UUID : pour Nova, une instance ou une migration}}

\newglossaryentry{trait}{name={trait}, sort={trait},
  description={Caractéristique qualitative d'un fournisseur (par exemple \texttt{HW\_CPU\_X86\_AVX2} ou \texttt{STORAGE\_DISK\_SSD}) ; les traits non standard commencent par \texttt{CUSTOM\_}}}

\newglossaryentry{claim}{name={réservation (claim)}, sort={reservation}, text={réservation},
  description={Création d'une allocation dans Placement pour une instance, faite par l'ordonnanceur une fois l'hôte choisi ; si elle échoue, c'est qu'un autre consommateur a pris les ressources entre-temps}}

% ---- Glance ----------------------------------------------------------
\newglossaryentry{image}{name={image}, sort={image},
  description={Fichier disque d'une machine virtuelle (système installé) qu'on duplique pour créer des instances ; Glance en conserve les données et les métadonnées}}

\newglossaryentry{qcow2}{name={qcow2}, sort={qcow2},
  description={Format de disque de QEMU à copie sur écriture : le fichier ne grossit qu'avec les données écrites et peut s'appuyer sur un fichier de base}}

\newglossaryentry{store}{name={store (backend de stockage)}, sort={store}, text={store},
  description={Système où Glance garde les données des images : système de fichiers, Ceph (rbd), Swift, Cinder, VMware ; en lecture seule pour http}}

% ---- Neutron ---------------------------------------------------------
\newglossaryentry{reseau}{name={réseau (network)}, sort={reseau}, text={réseau}, plural={réseaux},
  description={Dans Neutron, domaine de diffusion de niveau 2 isolé, auquel on rattache des ports}}

\newglossaryentry{sousreseau}{name={sous-réseau (subnet)}, sort={sous-reseau}, text={sous-réseau}, plural={sous-réseaux},
  description={Plage d'adresses IP associée à un réseau Neutron, avec sa passerelle et son serveur DHCP éventuel}}

\newglossaryentry{port}{name={port}, sort={port},
  description={Point de connexion d'un équipement (en général la carte réseau d'une instance) à un réseau Neutron ; il porte une adresse MAC, des adresses IP et des groupes de sécurité}}

\newglossaryentry{routeur}{name={routeur (router)}, sort={routeur}, text={routeur},
  description={Objet Neutron qui relie des réseaux entre eux et vers l'extérieur, avec traduction d'adresses (NAT)}}

\newglossaryentry{groupesecurite}{name={groupe de sécurité}, sort={groupe de securite}, text={groupe de sécurité}, plural={groupes de sécurité},
  description={Ensemble de règles de pare-feu appliqué aux ports : par défaut tout le trafic entrant est refusé et tout le trafic sortant autorisé ; les règles sont avec état}}

\newglossaryentry{ipflottante}{name={IP flottante (floating IP)}, sort={IP flottante}, text={IP flottante},
  description={Adresse publique associée à la demande à un port d'instance, et traduite en son adresse privée par le routeur}}

\newglossaryentry{provider}{name={réseau provider}, sort={reseau provider}, text={réseau provider},
  description={Réseau créé par l'administrateur et rattaché directement à un réseau physique existant (flat ou VLAN)}}

\newglossaryentry{selfservice}{name={réseau self-service}, sort={reseau self-service}, text={réseau self-service},
  description={Réseau virtuel (overlay) créé par un projet sans droits d'administrateur ; il a besoin d'un routeur pour sortir}}

\newglossaryentry{overlay}{name={overlay (réseau superposé)}, sort={overlay}, text={overlay},
  description={Réseau virtuel dont les trames sont encapsulées (VXLAN, Geneve) dans des paquets IP du réseau physique}}

\newglossaryentry{vlan}{name={VLAN}, sort={VLAN},
  description={Réseau local virtuel 802.1Q identifié par un numéro de 1 à 4094 ; la limite de 4094 réseaux restreint son usage pour les projets}}

\newglossaryentry{vxlan}{name={VXLAN, Geneve}, sort={VXLAN}, text={VXLAN},
  description={Protocoles de tunnel qui encapsulent des trames Ethernet dans des paquets UDP/IP ; ils servent aux réseaux overlay et n'ont pas la limite des VLAN}}

\newglossaryentry{ml2}{name={ML2 (Modular Layer 2)}, sort={ML2}, text={ML2},
  description={Plugin de base de Neutron : il orchestre des \emph{type drivers} (flat, VLAN, VXLAN\dots) et des \emph{mechanism drivers} (Open vSwitch, OVN\dots)}}

\newglossaryentry{ovs}{name={Open vSwitch (OVS)}, sort={Open vSwitch}, text={Open vSwitch},
  description={Commutateur virtuel programmable : il relie sur chaque nœud les interfaces des instances, les tunnels et les interfaces physiques, selon des règles (\emph{flows}) OpenFlow}}

\newglossaryentry{ovn}{name={OVN (Open Virtual Network)}, sort={OVN}, text={OVN},
  description={Système de réseau virtuel bâti sur Open vSwitch, avec deux bases de données et un démon par nœud ; il distribue le routage, le DHCP et les règles de sécurité}}

\newglossaryentry{agentneutron}{name={agent Neutron}, sort={agent Neutron}, text={agent},
  description={Processus qui, sur un nœud, traduit les objets Neutron en configuration réelle : agent L2 (Open vSwitch), L3 (routeurs), DHCP, métadonnées}}

\newglossaryentry{nat}{name={NAT}, sort={NAT},
  description={Traduction d'adresses réseau : le routeur remplace l'adresse source ou destination d'un paquet (SNAT pour la sortie, DNAT pour l'entrée)}}

% ---- Cinder -----------------------------------------------------------
\newglossaryentry{volume}{name={volume}, sort={volume},
  description={Disque virtuel persistant fourni par Cinder : il a une taille et un type, et peut être attaché à une instance, détaché, puis rattaché ailleurs sans perdre ses données}}

\newglossaryentry{snapshot}{name={snapshot}, sort={snapshot},
  description={Image ponctuelle d'un volume, conservée sur le même backend ; elle sert au retour arrière et au clonage, mais disparaît avec le backend}}

\newglossaryentry{sauvegarde}{name={sauvegarde (backup)}, sort={sauvegarde}, text={sauvegarde},
  description={Copie complète ou incrémentale d'un volume vers un dépôt distinct (Swift par défaut) ; elle survit à la perte du backend}}

\newglossaryentry{typevolume}{name={type de volume}, sort={type de volume}, plural={types de volume},
  description={Étiquette qui porte des propriétés (\emph{extra specs}) : backend visé, qualité de service, accès multiple ; Cinder s'en sert pour choisir le backend d'un volume}}

\newglossaryentry{backendcinder}{name={backend (Cinder)}, sort={backend}, text={backend}, plural={backends},
  description={Système de stockage qui héberge réellement les volumes (LVM, Ceph, NFS, baie SAN) ; \texttt{cinder-volume} le pilote par un driver}}

\newglossaryentry{qos}{name={QoS (qualité de service)}, sort={QoS}, text={QoS},
  description={Limite de débit appliquée à un volume, en opérations ou en octets par seconde ; elle s'associe à un type de volume}}

\newglossaryentry{osbrick}{name={os-brick}, sort={os-brick},
  description={Bibliothèque qui connecte un hôte à un volume (iSCSI, RBD, NFS\dots) ; Nova l'utilise sur le nœud de calcul pour attacher les volumes}}

% ---- Horizon ----------------------------------------------------------
\newglossaryentry{django}{name={Django}, sort={Django},
  description={Cadriciel Web en Python sur lequel est écrit Horizon ; il fournit les sessions, l'authentification par des backends et le système de gabarits}}

\newglossaryentry{dashboard}{name={dashboard (tableau de bord)}, sort={dashboard}, text={dashboard},
  description={Entrée de la navigation principale d'Horizon (\emph{Project}, \emph{Admin}, \emph{Identity}, \emph{Settings}) ; il regroupe des groupes de panneaux}}

\newglossaryentry{panneau}{name={panneau (panel)}, sort={panneau}, text={panneau}, plural={panneaux},
  description={Page de la navigation latérale d'un dashboard, déclarée dans un fichier \texttt{panel.py} ; elle a ses routes, ses vues et ses gabarits}}

\newglossaryentry{memcached}{name={memcached}, sort={memcached},
  description={Service de cache en mémoire, partagé entre processus et serveurs ; Horizon l'utilise couramment pour conserver les sessions}}

% ---- Octavia ----------------------------------------------------------
\newglossaryentry{loadbalancer}{name={répartiteur de charge (load balancer)}, sort={repartiteur de charge}, text={répartiteur de charge}, plural={répartiteurs de charge},
  description={Service qui reçoit le trafic sur une adresse unique (la VIP) et le distribue entre plusieurs serveurs ; c'est l'objet principal que l'utilisateur crée dans Octavia}}

\newglossaryentry{listener}{name={listener}, sort={listener}, plural={listeners},
  description={Point d'écoute d'un répartiteur : un protocole et un port (par exemple HTTP, 80) sur la VIP ; il envoie le trafic reçu vers un pool}}

\newglossaryentry{poolocta}{name={pool (Octavia)}, sort={pool}, text={pool}, plural={pools},
  description={Groupe de serveurs qui reçoivent le trafic d'un listener, avec un algorithme de répartition et, au besoin, un health monitor}}

\newglossaryentry{memberocta}{name={member (Octavia)}, sort={member}, text={member}, plural={members},
  description={Serveur d'un pool, repéré par une adresse IP et un port ; il peut être une instance Nova ou n'importe quelle adresse joignable}}

\newglossaryentry{healthmonitor}{name={health monitor}, sort={health monitor},
  description={Sonde que le répartiteur envoie à chaque member (TCP, HTTP, HTTPS, PING\dots) ; un member qui échoue est retiré de la distribution jusqu'à ce qu'il réponde de nouveau}}

\newglossaryentry{amphore}{name={amphore}, sort={amphore}, plural={amphores},
  description={Instance Nova, créée par Octavia, qui exécute HAProxy et porte la VIP ; c'est elle qui écoute et répartit réellement le trafic}}

\newglossaryentry{vip}{name={VIP (adresse virtuelle)}, sort={VIP}, text={VIP},
  description={Adresse IP du répartiteur de charge, celle que les clients utilisent ; elle est portée par un port Neutron et peut recevoir une adresse flottante}}

\newglossaryentry{l7}{name={politique L7 (L7 policy)}, sort={politique L7}, text={politique L7}, plural={politiques L7},
  description={Règle de routage au niveau applicatif (chemin d'URL, en-tête, nom d'hôte\dots) qui redirige une requête vers un autre pool ou vers une URL}}

\newglossaryentry{haproxy}{name={HAProxy}, sort={HAProxy},
  description={Logiciel libre de répartition de charge TCP et HTTP ; l'amphore de référence d'Octavia l'exécute}}

% ---- Swift ------------------------------------------------------------
\newglossaryentry{objetswift}{name={objet (Swift)}, sort={objet}, text={objet}, plural={objets},
  description={Donnée stockée dans Swift : le contenu d'un fichier et ses métadonnées, adressé par son chemin \texttt{/compte/conteneur/objet}}}

\newglossaryentry{conteneur}{name={conteneur (Swift)}, sort={conteneur}, text={conteneur}, plural={conteneurs},
  description={Espace de noms d'un compte Swift qui regroupe des objets ; il porte les droits d'accès, la politique de stockage et d'autres métadonnées}}

\newglossaryentry{proxyswift}{name={proxy (Swift)}, sort={proxy}, text={proxy}, plural={proxys},
  description={Service \texttt{swift-proxy-server} : seul point d'entrée de l'API, il authentifie, consulte le ring et relaie les requêtes vers les serveurs de stockage}}

\newglossaryentry{ring}{name={ring}, sort={ring}, plural={rings},
  description={Fichier qui associe chaque partition à des disques précis ; Swift en tient un pour les comptes, un pour les conteneurs et un par politique de stockage d'objets}}

\newglossaryentry{partition}{name={partition (Swift)}, sort={partition}, text={partition}, plural={partitions},
  description={Unité de répartition du ring : un intervalle des empreintes MD5 des chemins ; le ring place ses répliques sur des disques de zones différentes}}

\newglossaryentry{handoff}{name={handoff}, sort={handoff}, plural={handoffs},
  description={Disque de remplacement choisi par le ring lorsqu'un disque normal est en panne ; la réplication ramène ensuite les données à leur place}}

\newglossaryentry{erasurecoding}{name={codage d'effacement (erasure coding)}, sort={codage d'effacement}, text={codage d'effacement},
  description={Méthode de redondance qui découpe un objet en fragments de données et de parité ; Swift l'offre par une politique de stockage, pour un surcoût moindre que la triple réplication}}

\newglossaryentry{politiquestockage}{name={politique de stockage (storage policy)}, sort={politique de stockage}, text={politique de stockage}, plural={politiques de stockage},
  description={Choix, par conteneur, d'un ring d'objets et donc du nombre de répliques, du type de disques ou du codage d'effacement}}

% ---- Heat -------------------------------------------------------------
\newglossaryentry{modeleheat}{name={modèle (template)}, sort={modele}, text={modèle}, plural={modèles},
  description={Fichier qui décrit l'infrastructure voulue (ressources, paramètres, sorties) ; Heat le lit pour créer une stack. Son format natif est HOT, en YAML}}

\newglossaryentry{stack}{name={stack}, sort={stack}, plural={stacks},
  description={Ensemble de ressources créé par Heat à partir d'un modèle ; c'est l'unité qu'on crée, met à jour et supprime}}

\newglossaryentry{ressourceheat}{name={ressource (Heat)}, sort={ressource}, text={ressource}, plural={ressources},
  description={Élément d'infrastructure d'une stack (réseau, serveur, volume\dots), instancié d'après un type de la forme \texttt{OS::Service::Type}}}

\newglossaryentry{hot}{name={HOT (Heat Orchestration Template)}, sort={HOT}, text={HOT},
  description={Format de modèle natif de Heat, écrit en YAML ; sa version (\texttt{heat\_template\_version}) détermine les fonctions disponibles}}

\newglossaryentry{environnementheat}{name={environnement (Heat)}, sort={environnement}, text={environnement},
  description={Fichier passé avec \texttt{-e} qui fournit des valeurs de paramètres et peut remplacer des types de ressources (\texttt{resource\_registry})}}

\newglossaryentry{graphedeps}{name={graphe de dépendances}, sort={graphe de dependances}, text={graphe de dépendances},
  description={Graphe orienté que Heat construit à partir des références entre ressources ; il fixe l'ordre de création, et la suppression suit l'ordre inverse}}

% ---- Ironic -----------------------------------------------------------
\newglossaryentry{baremetal}{name={bare metal}, sort={bare metal}, text={bare metal},
  description={Serveur physique utilisé directement par un seul locataire, sans hyperviseur ; Ironic le provisionne et Nova le présente comme une instance}}

\newglossaryentry{noeudironic}{name={nœud (Ironic)}, sort={noeud (Ironic)}, text={nœud}, plural={nœuds},
  description={Enregistrement d'Ironic qui représente un serveur physique : son type de matériel, ses identifiants de gestion, ses ports, sa classe de ressources et son état de provisionnement}}

\newglossaryentry{portironic}{name={port (Ironic)}, sort={port (Ironic)}, text={port}, plural={ports},
  description={Carte réseau d'un nœud, déclarée dans Ironic par son adresse MAC ; à ne pas confondre avec le port Neutron, qui lui est lié lors du raccordement au réseau}}

\newglossaryentry{conductorironic}{name={conductor (Ironic)}, sort={conductor (Ironic)}, text={conductor}, plural={conductors},
  description={Composant d'Ironic (\texttt{ironic-conductor}) qui agit sur les serveurs physiques : alimentation, déploiement, nettoyage, inspection ; chaque conductor gère une partie des nœuds}}

\newglossaryentry{bmc}{name={BMC (contrôleur de gestion)}, sort={BMC}, text={BMC},
  description={\emph{Baseboard Management Controller} : contrôleur intégré au serveur, doté de sa propre adresse réseau ; il reçoit les ordres d'alimentation et de démarrage (IPMI, Redfish)}}

\newglossaryentry{ipmi}{name={IPMI}, sort={IPMI}, text={IPMI},
  description={\emph{Intelligent Platform Management Interface} : protocole ancien de gestion à distance par le BMC (alimentation, périphérique de démarrage) ; Ironic l'utilise avec la commande \texttt{ipmitool}}}

\newglossaryentry{redfish}{name={Redfish}, sort={Redfish}, text={Redfish},
  description={Protocole de gestion à distance des serveurs par API REST, proposé par le BMC ; Ironic l'utilise avec le type de matériel \texttt{redfish}}}

\newglossaryentry{pxe}{name={PXE (démarrage par le réseau)}, sort={PXE}, text={PXE},
  description={\emph{Preboot Execution Environment} : démarrage par le réseau, où un serveur DHCP indique au serveur le programme à charger et un serveur TFTP (ou HTTP) le fournit ; iPXE en est une variante}}

\newglossaryentry{ipa}{name={IPA (Ironic Python Agent)}, sort={IPA}, text={IPA},
  description={Agent Python qui s'exécute dans un ramdisk sur le serveur : il donne au conductor un accès à distance et réalise de l'intérieur l'écriture de l'image, le nettoyage et l'inspection}}

\newglossaryentry{hardwaretype}{name={type de matériel (hardware type)}, sort={type de materiel}, text={type de matériel}, plural={types de matériel},
  description={Ensemble d'interfaces matérielles qui définit comment Ironic pilote une famille de serveurs (\texttt{ipmi}, \texttt{redfish} ou un type propre à un constructeur) ; on parle aussi de pilote (\emph{driver})}}

\newglossaryentry{interfacemat}{name={interface matérielle (hardware interface)}, sort={interface materielle}, text={interface matérielle}, plural={interfaces matérielles},
  description={Brique d'un type de matériel qui gère un aspect du serveur : alimentation, démarrage, déploiement, réseau, inspection, RAID\dots}}

\newglossaryentry{nettoyage}{name={nettoyage (cleaning)}, sort={nettoyage}, text={nettoyage},
  description={Étapes qu'Ironic exécute sur un nœud entre deux locataires, ou avant le premier : effacement des disques, remise du matériel dans une configuration connue ; l'état \texttt{cleaning} les regroupe}}

\newglossaryentry{ramdisk}{name={ramdisk de déploiement}, sort={ramdisk}, text={ramdisk},
  description={Petit système Linux chargé en mémoire par démarrage réseau, qui exécute l'agent IPA pour déployer, nettoyer ou inspecter un nœud}}

\newglossaryentry{etatprov}{name={état de provisionnement (provision state)}, sort={etat de provisionnement}, text={état de provisionnement}, plural={états de provisionnement},
  description={Étape du cycle de vie d'un nœud dans la machine à états d'Ironic (\texttt{enroll}, \texttt{manageable}, \texttt{available}, \texttt{active}\dots)}}
)
%  Utilisation dans le texte : \gls{iaas} (le mot apparaît en rouge) ; \Gls{...} met la
%  première lettre en capitale.
%  Toutes les entrées sont imprimées en fin de document (\glsaddall dans main.tex).
%  La clé sort={...} donne le texte brut utilisé pour le tri (accents, commandes).
%  Compilation : deux passes de pdflatex suffisent (\makenoidxglossaries).
% =====================================================================

\newglossaryentry{apirest}{name={API REST}, sort={API REST},
  description={Interface d'accès à un service par requêtes HTTP ; chaque service OpenStack expose la sienne, et la ligne de commande comme Horizon n'en sont que des clients}}

\newglossaryentry{appcred}{name={application credential}, sort={application credential},
  description={Identifiant et secret délivrés à une application pour qu'elle s'authentifie à la place de l'utilisateur, avec seulement une partie de ses rôles et, éventuellement, une date d'expiration}}

\newglossaryentry{authentification}{name={authentification}, sort={authentification},
  description={Processus qui confirme l'identité d'un utilisateur ; Keystone valide les identifiants puis délivre un token}}

\newglossaryentry{autorisation}{name={autorisation}, sort={autorisation},
  description={Décision d'autoriser ou de refuser une action à un utilisateur authentifié, d'après ses rôles ; chaque service OpenStack l'applique avec ses propres règles}}

\newglossaryentry{cadf}{name={CADF}, sort={CADF},
  description={\emph{Cloud Auditing Data Federation} : format normalisé d'événements d'audit, utilisé par défaut pour les notifications de Keystone}}

\newglossaryentry{catalogue}{name={catalogue de services}, sort={catalogue de services}, text={catalogue},
  description={Liste des services de la plateforme et de leurs endpoints, par région et par interface ; Keystone l'inclut dans les tokens limités à un projet}}

\newglossaryentry{domaine}{name={domaine}, sort={domaine},
  description={Conteneur de plus haut niveau de Keystone : il regroupe des projets, des utilisateurs et des groupes qui forment une même frontière administrative}}

\newglossaryentry{endpoint}{name={endpoint}, sort={endpoint},
  description={Adresse réseau, généralement une URL, à laquelle on accède à un service ; Keystone en distingue trois interfaces : \texttt{public}, \texttt{internal} et \texttt{admin}}}

\newglossaryentry{federation}{name={fédération d'identité}, sort={federation d'identite},
  description={Délégation de l'authentification à un fournisseur d'identité externe de confiance, par SAML 2.0 ou OpenID Connect}}

\newglossaryentry{fernet}{name={Fernet}, sort={Fernet},
  description={Format de token par défaut de Keystone : chaîne opaque chiffrée (AES-128-CBC) et signée (HMAC-SHA256) avec des clés conservées dans un dépôt local}}

\newglossaryentry{groupe}{name={groupe}, sort={groupe},
  description={Ensemble d'utilisateurs d'un domaine ; un rôle attribué à un groupe s'applique à tous ses membres}}

\newglossaryentry{hyperviseur}{name={hyperviseur}, sort={hyperviseur},
  description={Logiciel qui fait tourner plusieurs machines virtuelles isolées sur un même serveur physique (par exemple KVM) ; OpenStack le pilote sans le remplacer}}

\newglossaryentry{iaas}{name={IaaS}, sort={IaaS},
  description={\emph{Infrastructure as a Service} : modèle de cloud où l'infrastructure (calcul, réseau, stockage) est fournie à la demande ; l'utilisateur gère ce qui s'exécute dessus}}

\newglossaryentry{idp}{name={fournisseur d'identité (IdP)}, sort={fournisseur d'identite}, text={fournisseur d'identité}, plural={fournisseurs d'identité},
  description={Service externe qui authentifie l'utilisateur et transmet à Keystone une assertion signée contenant son identité et ses attributs}}

\newglossaryentry{jws}{name={JWS}, sort={JWS},
  description={\emph{JSON Web Signature} : second format de token de Keystone, signé avec une clé asymétrique (ES256) ; son contenu n'est pas opaque}}

\newglossaryentry{ldap}{name={LDAP}, sort={LDAP},
  description={\emph{Lightweight Directory Access Protocol} : protocole d'accès à un annuaire (par exemple Active Directory) ; Keystone peut y lire utilisateurs et groupes, en lecture seule}}

\newglossaryentry{mapping}{name={mapping}, sort={mapping},
  description={Ensemble de règles qui traduisent les attributs fournis par un fournisseur d'identité en utilisateur, groupes et projets Keystone}}

\newglossaryentry{mfa}{name={authentification multifacteur (MFA)}, sort={authentification multifacteur}, text={authentification multifacteur},
  description={Authentification qui exige plusieurs méthodes successives, par exemple un mot de passe puis un code TOTP}}

\newglossaryentry{middleware}{name={middleware}, sort={middleware},
  description={Composant WSGI placé devant l'application d'un service ; \texttt{keystonemiddleware} (module \texttt{auth\_token}) y valide le token de chaque requête auprès de Keystone}}

\newglossaryentry{oidc}{name={OpenID Connect (OIDC)}, sort={OpenID Connect}, text={OpenID Connect},
  description={Protocole de fédération au format JSON, construit sur OAuth 2.0}}

\newglossaryentry{policy}{name={policy}, sort={policy}, plural={policies},
  description={Règle d'accès évaluée par chaque service (bibliothèque \texttt{oslo.policy}) pour décider si les rôles présents dans le token autorisent l'action demandée}}

\newglossaryentry{portee}{name={portée (scope)}, sort={portee}, text={portée}, plural={portées},
  description={Périmètre sur lequel un token confère des droits : un projet, un domaine ou le système entier ; un token non délimité n'a aucune portée}}

\newglossaryentry{projet}{name={projet (tenant)}, sort={projet}, text={projet}, plural={projets},
  description={Conteneur qui regroupe ou isole des ressources ; « tenant » est l'ancien nom du projet}}

\newglossaryentry{rbac}{name={RBAC}, sort={RBAC},
  description={\emph{Role-Based Access Control} : contrôle d'accès fondé sur les rôles attribués aux utilisateurs et aux groupes}}

\newglossaryentry{region}{name={région}, sort={region},
  description={Division générale d'un déploiement OpenStack, qui peut contenir des sous-régions ; chaque endpoint du catalogue est rattaché à une région}}

\newglossaryentry{role}{name={rôle}, sort={role},
  description={Ensemble de droits attribué à un utilisateur ou à un groupe sur un projet, un domaine ou le système ; les services l'interprètent pour décider des accès}}

\newglossaryentry{saml2}{name={SAML 2.0}, sort={SAML},
  description={\emph{Security Assertion Markup Language} : protocole de fédération au format XML, courant dans les universités et les entreprises}}

\newglossaryentry{sp}{name={fournisseur de services (SP)}, sort={fournisseur de services}, text={fournisseur de services}, plural={fournisseurs de services},
  description={Dans une fédération, service qui s'appuie sur un fournisseur d'identité ; ici Keystone, dont les tokens servent aux autres services OpenStack}}

\newglossaryentry{sso}{name={SSO}, sort={SSO},
  description={\emph{Single Sign-On}, authentification unique : une seule connexion donne accès à plusieurs services}}

\newglossaryentry{token}{name={token (jeton)}, sort={token}, text={token}, plural={tokens},
  description={Chaîne de caractères qui prouve l'identité d'un utilisateur et donne accès aux API ; valide pour une durée limitée, elle peut être révoquée à tout moment}}

\newglossaryentry{totp}{name={TOTP}, sort={TOTP},
  description={\emph{Time-based One-Time Password} : code à usage unique de six chiffres, calculé à partir d'un secret partagé et de l'heure ; second facteur pris en charge par Keystone}}

\newglossaryentry{trust}{name={trust}, sort={trust},
  description={Délégation par un utilisateur (le \emph{trustor}) d'une partie de ses rôles sur un projet à un autre utilisateur (le \emph{trustee})}}

\newglossaryentry{utilisateur}{name={utilisateur}, sort={utilisateur},
  description={Représentation numérique d'une personne, d'un système ou d'un service qui utilise les services OpenStack}}

\newglossaryentry{wsgi}{name={WSGI}, sort={WSGI},
  description={\emph{Web Server Gateway Interface} : interface standard entre un serveur web et une application Python ; Keystone s'exécute derrière Apache grâce à \texttt{mod\_wsgi}}}

% ---- Nova et Placement ------------------------------------------------
\newglossaryentry{instance}{name={instance}, sort={instance},
  description={Machine virtuelle créée et gérée par Nova ; on parle aussi de serveur (\texttt{openstack server})}}

\newglossaryentry{flavor}{name={gabarit (flavor)}, sort={gabarit}, text={gabarit},
  description={Modèle de taille d'une instance : nombre de vCPU, mémoire, disque racine, disque éphémère, swap, et propriétés supplémentaires (\emph{extra specs})}}

\newglossaryentry{noeudcalcul}{name={nœud de calcul (compute node)}, sort={noeud de calcul}, text={nœud de calcul},
  description={Serveur physique qui héberge les instances : il exécute \texttt{nova-compute} et un hyperviseur}}

\newglossaryentry{ordonnanceur}{name={ordonnanceur (scheduler)}, sort={ordonnanceur}, text={ordonnanceur},
  description={Composant de Nova (\texttt{nova-scheduler}) qui choisit le nœud de calcul où placer une instance}}

\newglossaryentry{conductor}{name={conductor}, sort={conductor},
  description={Composant de Nova (\texttt{nova-conductor}) qui coordonne les opérations longues (construction, redimensionnement, migration) et sert d'intermédiaire entre les nœuds de calcul et la base de données}}

\newglossaryentry{cellule}{name={cellule (cell)}, sort={cellule}, text={cellule},
  description={Groupe de nœuds de calcul de Nova qui possède sa propre base de données et sa propre file de messages ; tout déploiement a au moins une cellule}}

\newglossaryentry{libvirt}{name={libvirt}, sort={libvirt},
  description={Bibliothèque et service de pilotage des hyperviseurs (KVM/QEMU en particulier) que \texttt{nova-compute} utilise pour créer et contrôler les machines virtuelles}}

\newglossaryentry{busmessages}{name={bus de messages}, sort={bus de messages}, text={bus de messages},
  description={File de messages (RabbitMQ en général) par laquelle les composants d'un service échangent des appels distants (RPC, via \texttt{oslo.messaging})}}

\newglossaryentry{microversion}{name={microversion}, sort={microversion},
  description={Numéro de version fin d'une API OpenStack (pour Nova : 2.1, 2.2, 2.3 et ainsi de suite) : le client demande la version dont il a besoin, ce qui garde l'API compatible}}

\newglossaryentry{migrationchaud}{name={migration à chaud (live migration)}, sort={migration a chaud}, text={migration à chaud},
  description={Déplacement d'une instance en cours d'exécution d'un nœud de calcul vers un autre, sans l'arrêter}}

\newglossaryentry{surallocation}{name={surallocation (overcommit)}, sort={surallocation}, text={surallocation},
  description={Attribution de plus de ressources virtuelles qu'il n'y en a de physiques ; elle se règle par un ratio d'allocation}}

\newglossaryentry{agregat}{name={agrégat d'hôtes (host aggregate)}, sort={agregat}, text={agrégat d'hôtes},
  description={Groupe de nœuds de calcul défini par l'administrateur, avec des métadonnées ; il sert à l'ordonnancement et à définir des zones de disponibilité}}

\newglossaryentry{zonedispo}{name={zone de disponibilité}, sort={zone de disponibilite}, text={zone de disponibilité},
  description={Groupe de nœuds de calcul que l'utilisateur peut viser à la création d'une instance (par exemple une salle ou une baie) ; c'est un agrégat d'hôtes visible des utilisateurs}}

\newglossaryentry{fournisseur}{name={fournisseur de ressources (resource provider)}, sort={fournisseur de ressources}, text={fournisseur de ressources},
  description={Entité de Placement qui offre un inventaire de ressources consommables, par exemple un nœud de calcul ou un pool de stockage partagé}}

\newglossaryentry{inventaire}{name={inventaire}, sort={inventaire},
  description={Quantité d'une classe de ressources offerte par un fournisseur : total, réservé, ratio d'allocation, taille minimale, maximale et pas d'allocation}}

\newglossaryentry{classeressources}{name={classe de ressources}, sort={classe de ressources}, text={classe de ressources},
  description={Type de ressource quantitative suivie par Placement : \texttt{VCPU}, \texttt{MEMORY\_MB}, \texttt{DISK\_GB}, ou une classe personnalisée préfixée par \texttt{CUSTOM\_}}}

\newglossaryentry{allocation}{name={allocation}, sort={allocation},
  description={Quantité de ressources d'un fournisseur attribuée à un consommateur ; elle diminue d'autant la capacité disponible}}

\newglossaryentry{consommateur}{name={consommateur (consumer)}, sort={consommateur}, text={consommateur},
  description={Entité qui détient des allocations dans Placement, identifiée par un UUID : pour Nova, une instance ou une migration}}

\newglossaryentry{trait}{name={trait}, sort={trait},
  description={Caractéristique qualitative d'un fournisseur (par exemple \texttt{HW\_CPU\_X86\_AVX2} ou \texttt{STORAGE\_DISK\_SSD}) ; les traits non standard commencent par \texttt{CUSTOM\_}}}

\newglossaryentry{claim}{name={réservation (claim)}, sort={reservation}, text={réservation},
  description={Création d'une allocation dans Placement pour une instance, faite par l'ordonnanceur une fois l'hôte choisi ; si elle échoue, c'est qu'un autre consommateur a pris les ressources entre-temps}}

% ---- Glance ----------------------------------------------------------
\newglossaryentry{image}{name={image}, sort={image},
  description={Fichier disque d'une machine virtuelle (système installé) qu'on duplique pour créer des instances ; Glance en conserve les données et les métadonnées}}

\newglossaryentry{qcow2}{name={qcow2}, sort={qcow2},
  description={Format de disque de QEMU à copie sur écriture : le fichier ne grossit qu'avec les données écrites et peut s'appuyer sur un fichier de base}}

\newglossaryentry{store}{name={store (backend de stockage)}, sort={store}, text={store},
  description={Système où Glance garde les données des images : système de fichiers, Ceph (rbd), Swift, Cinder, VMware ; en lecture seule pour http}}

% ---- Neutron ---------------------------------------------------------
\newglossaryentry{reseau}{name={réseau (network)}, sort={reseau}, text={réseau}, plural={réseaux},
  description={Dans Neutron, domaine de diffusion de niveau 2 isolé, auquel on rattache des ports}}

\newglossaryentry{sousreseau}{name={sous-réseau (subnet)}, sort={sous-reseau}, text={sous-réseau}, plural={sous-réseaux},
  description={Plage d'adresses IP associée à un réseau Neutron, avec sa passerelle et son serveur DHCP éventuel}}

\newglossaryentry{port}{name={port}, sort={port},
  description={Point de connexion d'un équipement (en général la carte réseau d'une instance) à un réseau Neutron ; il porte une adresse MAC, des adresses IP et des groupes de sécurité}}

\newglossaryentry{routeur}{name={routeur (router)}, sort={routeur}, text={routeur},
  description={Objet Neutron qui relie des réseaux entre eux et vers l'extérieur, avec traduction d'adresses (NAT)}}

\newglossaryentry{groupesecurite}{name={groupe de sécurité}, sort={groupe de securite}, text={groupe de sécurité}, plural={groupes de sécurité},
  description={Ensemble de règles de pare-feu appliqué aux ports : par défaut tout le trafic entrant est refusé et tout le trafic sortant autorisé ; les règles sont avec état}}

\newglossaryentry{ipflottante}{name={IP flottante (floating IP)}, sort={IP flottante}, text={IP flottante},
  description={Adresse publique associée à la demande à un port d'instance, et traduite en son adresse privée par le routeur}}

\newglossaryentry{provider}{name={réseau provider}, sort={reseau provider}, text={réseau provider},
  description={Réseau créé par l'administrateur et rattaché directement à un réseau physique existant (flat ou VLAN)}}

\newglossaryentry{selfservice}{name={réseau self-service}, sort={reseau self-service}, text={réseau self-service},
  description={Réseau virtuel (overlay) créé par un projet sans droits d'administrateur ; il a besoin d'un routeur pour sortir}}

\newglossaryentry{overlay}{name={overlay (réseau superposé)}, sort={overlay}, text={overlay},
  description={Réseau virtuel dont les trames sont encapsulées (VXLAN, Geneve) dans des paquets IP du réseau physique}}

\newglossaryentry{vlan}{name={VLAN}, sort={VLAN},
  description={Réseau local virtuel 802.1Q identifié par un numéro de 1 à 4094 ; la limite de 4094 réseaux restreint son usage pour les projets}}

\newglossaryentry{vxlan}{name={VXLAN, Geneve}, sort={VXLAN}, text={VXLAN},
  description={Protocoles de tunnel qui encapsulent des trames Ethernet dans des paquets UDP/IP ; ils servent aux réseaux overlay et n'ont pas la limite des VLAN}}

\newglossaryentry{ml2}{name={ML2 (Modular Layer 2)}, sort={ML2}, text={ML2},
  description={Plugin de base de Neutron : il orchestre des \emph{type drivers} (flat, VLAN, VXLAN\dots) et des \emph{mechanism drivers} (Open vSwitch, OVN\dots)}}

\newglossaryentry{ovs}{name={Open vSwitch (OVS)}, sort={Open vSwitch}, text={Open vSwitch},
  description={Commutateur virtuel programmable : il relie sur chaque nœud les interfaces des instances, les tunnels et les interfaces physiques, selon des règles (\emph{flows}) OpenFlow}}

\newglossaryentry{ovn}{name={OVN (Open Virtual Network)}, sort={OVN}, text={OVN},
  description={Système de réseau virtuel bâti sur Open vSwitch, avec deux bases de données et un démon par nœud ; il distribue le routage, le DHCP et les règles de sécurité}}

\newglossaryentry{agentneutron}{name={agent Neutron}, sort={agent Neutron}, text={agent},
  description={Processus qui, sur un nœud, traduit les objets Neutron en configuration réelle : agent L2 (Open vSwitch), L3 (routeurs), DHCP, métadonnées}}

\newglossaryentry{nat}{name={NAT}, sort={NAT},
  description={Traduction d'adresses réseau : le routeur remplace l'adresse source ou destination d'un paquet (SNAT pour la sortie, DNAT pour l'entrée)}}

% ---- Cinder -----------------------------------------------------------
\newglossaryentry{volume}{name={volume}, sort={volume},
  description={Disque virtuel persistant fourni par Cinder : il a une taille et un type, et peut être attaché à une instance, détaché, puis rattaché ailleurs sans perdre ses données}}

\newglossaryentry{snapshot}{name={snapshot}, sort={snapshot},
  description={Image ponctuelle d'un volume, conservée sur le même backend ; elle sert au retour arrière et au clonage, mais disparaît avec le backend}}

\newglossaryentry{sauvegarde}{name={sauvegarde (backup)}, sort={sauvegarde}, text={sauvegarde},
  description={Copie complète ou incrémentale d'un volume vers un dépôt distinct (Swift par défaut) ; elle survit à la perte du backend}}

\newglossaryentry{typevolume}{name={type de volume}, sort={type de volume}, plural={types de volume},
  description={Étiquette qui porte des propriétés (\emph{extra specs}) : backend visé, qualité de service, accès multiple ; Cinder s'en sert pour choisir le backend d'un volume}}

\newglossaryentry{backendcinder}{name={backend (Cinder)}, sort={backend}, text={backend}, plural={backends},
  description={Système de stockage qui héberge réellement les volumes (LVM, Ceph, NFS, baie SAN) ; \texttt{cinder-volume} le pilote par un driver}}

\newglossaryentry{qos}{name={QoS (qualité de service)}, sort={QoS}, text={QoS},
  description={Limite de débit appliquée à un volume, en opérations ou en octets par seconde ; elle s'associe à un type de volume}}

\newglossaryentry{osbrick}{name={os-brick}, sort={os-brick},
  description={Bibliothèque qui connecte un hôte à un volume (iSCSI, RBD, NFS\dots) ; Nova l'utilise sur le nœud de calcul pour attacher les volumes}}

% ---- Horizon ----------------------------------------------------------
\newglossaryentry{django}{name={Django}, sort={Django},
  description={Cadriciel Web en Python sur lequel est écrit Horizon ; il fournit les sessions, l'authentification par des backends et le système de gabarits}}

\newglossaryentry{dashboard}{name={dashboard (tableau de bord)}, sort={dashboard}, text={dashboard},
  description={Entrée de la navigation principale d'Horizon (\emph{Project}, \emph{Admin}, \emph{Identity}, \emph{Settings}) ; il regroupe des groupes de panneaux}}

\newglossaryentry{panneau}{name={panneau (panel)}, sort={panneau}, text={panneau}, plural={panneaux},
  description={Page de la navigation latérale d'un dashboard, déclarée dans un fichier \texttt{panel.py} ; elle a ses routes, ses vues et ses gabarits}}

\newglossaryentry{memcached}{name={memcached}, sort={memcached},
  description={Service de cache en mémoire, partagé entre processus et serveurs ; Horizon l'utilise couramment pour conserver les sessions}}

% ---- Octavia ----------------------------------------------------------
\newglossaryentry{loadbalancer}{name={répartiteur de charge (load balancer)}, sort={repartiteur de charge}, text={répartiteur de charge}, plural={répartiteurs de charge},
  description={Service qui reçoit le trafic sur une adresse unique (la VIP) et le distribue entre plusieurs serveurs ; c'est l'objet principal que l'utilisateur crée dans Octavia}}

\newglossaryentry{listener}{name={listener}, sort={listener}, plural={listeners},
  description={Point d'écoute d'un répartiteur : un protocole et un port (par exemple HTTP, 80) sur la VIP ; il envoie le trafic reçu vers un pool}}

\newglossaryentry{poolocta}{name={pool (Octavia)}, sort={pool}, text={pool}, plural={pools},
  description={Groupe de serveurs qui reçoivent le trafic d'un listener, avec un algorithme de répartition et, au besoin, un health monitor}}

\newglossaryentry{memberocta}{name={member (Octavia)}, sort={member}, text={member}, plural={members},
  description={Serveur d'un pool, repéré par une adresse IP et un port ; il peut être une instance Nova ou n'importe quelle adresse joignable}}

\newglossaryentry{healthmonitor}{name={health monitor}, sort={health monitor},
  description={Sonde que le répartiteur envoie à chaque member (TCP, HTTP, HTTPS, PING\dots) ; un member qui échoue est retiré de la distribution jusqu'à ce qu'il réponde de nouveau}}

\newglossaryentry{amphore}{name={amphore}, sort={amphore}, plural={amphores},
  description={Instance Nova, créée par Octavia, qui exécute HAProxy et porte la VIP ; c'est elle qui écoute et répartit réellement le trafic}}

\newglossaryentry{vip}{name={VIP (adresse virtuelle)}, sort={VIP}, text={VIP},
  description={Adresse IP du répartiteur de charge, celle que les clients utilisent ; elle est portée par un port Neutron et peut recevoir une adresse flottante}}

\newglossaryentry{l7}{name={politique L7 (L7 policy)}, sort={politique L7}, text={politique L7}, plural={politiques L7},
  description={Règle de routage au niveau applicatif (chemin d'URL, en-tête, nom d'hôte\dots) qui redirige une requête vers un autre pool ou vers une URL}}

\newglossaryentry{haproxy}{name={HAProxy}, sort={HAProxy},
  description={Logiciel libre de répartition de charge TCP et HTTP ; l'amphore de référence d'Octavia l'exécute}}

% ---- Swift ------------------------------------------------------------
\newglossaryentry{objetswift}{name={objet (Swift)}, sort={objet}, text={objet}, plural={objets},
  description={Donnée stockée dans Swift : le contenu d'un fichier et ses métadonnées, adressé par son chemin \texttt{/compte/conteneur/objet}}}

\newglossaryentry{conteneur}{name={conteneur (Swift)}, sort={conteneur}, text={conteneur}, plural={conteneurs},
  description={Espace de noms d'un compte Swift qui regroupe des objets ; il porte les droits d'accès, la politique de stockage et d'autres métadonnées}}

\newglossaryentry{proxyswift}{name={proxy (Swift)}, sort={proxy}, text={proxy}, plural={proxys},
  description={Service \texttt{swift-proxy-server} : seul point d'entrée de l'API, il authentifie, consulte le ring et relaie les requêtes vers les serveurs de stockage}}

\newglossaryentry{ring}{name={ring}, sort={ring}, plural={rings},
  description={Fichier qui associe chaque partition à des disques précis ; Swift en tient un pour les comptes, un pour les conteneurs et un par politique de stockage d'objets}}

\newglossaryentry{partition}{name={partition (Swift)}, sort={partition}, text={partition}, plural={partitions},
  description={Unité de répartition du ring : un intervalle des empreintes MD5 des chemins ; le ring place ses répliques sur des disques de zones différentes}}

\newglossaryentry{handoff}{name={handoff}, sort={handoff}, plural={handoffs},
  description={Disque de remplacement choisi par le ring lorsqu'un disque normal est en panne ; la réplication ramène ensuite les données à leur place}}

\newglossaryentry{erasurecoding}{name={codage d'effacement (erasure coding)}, sort={codage d'effacement}, text={codage d'effacement},
  description={Méthode de redondance qui découpe un objet en fragments de données et de parité ; Swift l'offre par une politique de stockage, pour un surcoût moindre que la triple réplication}}

\newglossaryentry{politiquestockage}{name={politique de stockage (storage policy)}, sort={politique de stockage}, text={politique de stockage}, plural={politiques de stockage},
  description={Choix, par conteneur, d'un ring d'objets et donc du nombre de répliques, du type de disques ou du codage d'effacement}}

% ---- Heat -------------------------------------------------------------
\newglossaryentry{modeleheat}{name={modèle (template)}, sort={modele}, text={modèle}, plural={modèles},
  description={Fichier qui décrit l'infrastructure voulue (ressources, paramètres, sorties) ; Heat le lit pour créer une stack. Son format natif est HOT, en YAML}}

\newglossaryentry{stack}{name={stack}, sort={stack}, plural={stacks},
  description={Ensemble de ressources créé par Heat à partir d'un modèle ; c'est l'unité qu'on crée, met à jour et supprime}}

\newglossaryentry{ressourceheat}{name={ressource (Heat)}, sort={ressource}, text={ressource}, plural={ressources},
  description={Élément d'infrastructure d'une stack (réseau, serveur, volume\dots), instancié d'après un type de la forme \texttt{OS::Service::Type}}}

\newglossaryentry{hot}{name={HOT (Heat Orchestration Template)}, sort={HOT}, text={HOT},
  description={Format de modèle natif de Heat, écrit en YAML ; sa version (\texttt{heat\_template\_version}) détermine les fonctions disponibles}}

\newglossaryentry{environnementheat}{name={environnement (Heat)}, sort={environnement}, text={environnement},
  description={Fichier passé avec \texttt{-e} qui fournit des valeurs de paramètres et peut remplacer des types de ressources (\texttt{resource\_registry})}}

\newglossaryentry{graphedeps}{name={graphe de dépendances}, sort={graphe de dependances}, text={graphe de dépendances},
  description={Graphe orienté que Heat construit à partir des références entre ressources ; il fixe l'ordre de création, et la suppression suit l'ordre inverse}}

% ---- Ironic -----------------------------------------------------------
\newglossaryentry{baremetal}{name={bare metal}, sort={bare metal}, text={bare metal},
  description={Serveur physique utilisé directement par un seul locataire, sans hyperviseur ; Ironic le provisionne et Nova le présente comme une instance}}

\newglossaryentry{noeudironic}{name={nœud (Ironic)}, sort={noeud (Ironic)}, text={nœud}, plural={nœuds},
  description={Enregistrement d'Ironic qui représente un serveur physique : son type de matériel, ses identifiants de gestion, ses ports, sa classe de ressources et son état de provisionnement}}

\newglossaryentry{portironic}{name={port (Ironic)}, sort={port (Ironic)}, text={port}, plural={ports},
  description={Carte réseau d'un nœud, déclarée dans Ironic par son adresse MAC ; à ne pas confondre avec le port Neutron, qui lui est lié lors du raccordement au réseau}}

\newglossaryentry{conductorironic}{name={conductor (Ironic)}, sort={conductor (Ironic)}, text={conductor}, plural={conductors},
  description={Composant d'Ironic (\texttt{ironic-conductor}) qui agit sur les serveurs physiques : alimentation, déploiement, nettoyage, inspection ; chaque conductor gère une partie des nœuds}}

\newglossaryentry{bmc}{name={BMC (contrôleur de gestion)}, sort={BMC}, text={BMC},
  description={\emph{Baseboard Management Controller} : contrôleur intégré au serveur, doté de sa propre adresse réseau ; il reçoit les ordres d'alimentation et de démarrage (IPMI, Redfish)}}

\newglossaryentry{ipmi}{name={IPMI}, sort={IPMI}, text={IPMI},
  description={\emph{Intelligent Platform Management Interface} : protocole ancien de gestion à distance par le BMC (alimentation, périphérique de démarrage) ; Ironic l'utilise avec la commande \texttt{ipmitool}}}

\newglossaryentry{redfish}{name={Redfish}, sort={Redfish}, text={Redfish},
  description={Protocole de gestion à distance des serveurs par API REST, proposé par le BMC ; Ironic l'utilise avec le type de matériel \texttt{redfish}}}

\newglossaryentry{pxe}{name={PXE (démarrage par le réseau)}, sort={PXE}, text={PXE},
  description={\emph{Preboot Execution Environment} : démarrage par le réseau, où un serveur DHCP indique au serveur le programme à charger et un serveur TFTP (ou HTTP) le fournit ; iPXE en est une variante}}

\newglossaryentry{ipa}{name={IPA (Ironic Python Agent)}, sort={IPA}, text={IPA},
  description={Agent Python qui s'exécute dans un ramdisk sur le serveur : il donne au conductor un accès à distance et réalise de l'intérieur l'écriture de l'image, le nettoyage et l'inspection}}

\newglossaryentry{hardwaretype}{name={type de matériel (hardware type)}, sort={type de materiel}, text={type de matériel}, plural={types de matériel},
  description={Ensemble d'interfaces matérielles qui définit comment Ironic pilote une famille de serveurs (\texttt{ipmi}, \texttt{redfish} ou un type propre à un constructeur) ; on parle aussi de pilote (\emph{driver})}}

\newglossaryentry{interfacemat}{name={interface matérielle (hardware interface)}, sort={interface materielle}, text={interface matérielle}, plural={interfaces matérielles},
  description={Brique d'un type de matériel qui gère un aspect du serveur : alimentation, démarrage, déploiement, réseau, inspection, RAID\dots}}

\newglossaryentry{nettoyage}{name={nettoyage (cleaning)}, sort={nettoyage}, text={nettoyage},
  description={Étapes qu'Ironic exécute sur un nœud entre deux locataires, ou avant le premier : effacement des disques, remise du matériel dans une configuration connue ; l'état \texttt{cleaning} les regroupe}}

\newglossaryentry{ramdisk}{name={ramdisk de déploiement}, sort={ramdisk}, text={ramdisk},
  description={Petit système Linux chargé en mémoire par démarrage réseau, qui exécute l'agent IPA pour déployer, nettoyer ou inspecter un nœud}}

\newglossaryentry{etatprov}{name={état de provisionnement (provision state)}, sort={etat de provisionnement}, text={état de provisionnement}, plural={états de provisionnement},
  description={Étape du cycle de vie d'un nœud dans la machine à états d'Ironic (\texttt{enroll}, \texttt{manageable}, \texttt{available}, \texttt{active}\dots)}}
)
%  Utilisation dans le texte : \gls{iaas} (le mot apparaît en rouge) ; \Gls{...} met la
%  première lettre en capitale.
%  Toutes les entrées sont imprimées en fin de document (\glsaddall dans main.tex).
%  La clé sort={...} donne le texte brut utilisé pour le tri (accents, commandes).
%  Compilation : deux passes de pdflatex suffisent (\makenoidxglossaries).
% =====================================================================

\newglossaryentry{apirest}{name={API REST}, sort={API REST},
  description={Interface d'accès à un service par requêtes HTTP ; chaque service OpenStack expose la sienne, et la ligne de commande comme Horizon n'en sont que des clients}}

\newglossaryentry{appcred}{name={application credential}, sort={application credential},
  description={Identifiant et secret délivrés à une application pour qu'elle s'authentifie à la place de l'utilisateur, avec seulement une partie de ses rôles et, éventuellement, une date d'expiration}}

\newglossaryentry{authentification}{name={authentification}, sort={authentification},
  description={Processus qui confirme l'identité d'un utilisateur ; Keystone valide les identifiants puis délivre un token}}

\newglossaryentry{autorisation}{name={autorisation}, sort={autorisation},
  description={Décision d'autoriser ou de refuser une action à un utilisateur authentifié, d'après ses rôles ; chaque service OpenStack l'applique avec ses propres règles}}

\newglossaryentry{cadf}{name={CADF}, sort={CADF},
  description={\emph{Cloud Auditing Data Federation} : format normalisé d'événements d'audit, utilisé par défaut pour les notifications de Keystone}}

\newglossaryentry{catalogue}{name={catalogue de services}, sort={catalogue de services}, text={catalogue},
  description={Liste des services de la plateforme et de leurs endpoints, par région et par interface ; Keystone l'inclut dans les tokens limités à un projet}}

\newglossaryentry{domaine}{name={domaine}, sort={domaine},
  description={Conteneur de plus haut niveau de Keystone : il regroupe des projets, des utilisateurs et des groupes qui forment une même frontière administrative}}

\newglossaryentry{endpoint}{name={endpoint}, sort={endpoint},
  description={Adresse réseau, généralement une URL, à laquelle on accède à un service ; Keystone en distingue trois interfaces : \texttt{public}, \texttt{internal} et \texttt{admin}}}

\newglossaryentry{federation}{name={fédération d'identité}, sort={federation d'identite},
  description={Délégation de l'authentification à un fournisseur d'identité externe de confiance, par SAML 2.0 ou OpenID Connect}}

\newglossaryentry{fernet}{name={Fernet}, sort={Fernet},
  description={Format de token par défaut de Keystone : chaîne opaque chiffrée (AES-128-CBC) et signée (HMAC-SHA256) avec des clés conservées dans un dépôt local}}

\newglossaryentry{groupe}{name={groupe}, sort={groupe},
  description={Ensemble d'utilisateurs d'un domaine ; un rôle attribué à un groupe s'applique à tous ses membres}}

\newglossaryentry{hyperviseur}{name={hyperviseur}, sort={hyperviseur},
  description={Logiciel qui fait tourner plusieurs machines virtuelles isolées sur un même serveur physique (par exemple KVM) ; OpenStack le pilote sans le remplacer}}

\newglossaryentry{iaas}{name={IaaS}, sort={IaaS},
  description={\emph{Infrastructure as a Service} : modèle de cloud où l'infrastructure (calcul, réseau, stockage) est fournie à la demande ; l'utilisateur gère ce qui s'exécute dessus}}

\newglossaryentry{idp}{name={fournisseur d'identité (IdP)}, sort={fournisseur d'identite}, text={fournisseur d'identité}, plural={fournisseurs d'identité},
  description={Service externe qui authentifie l'utilisateur et transmet à Keystone une assertion signée contenant son identité et ses attributs}}

\newglossaryentry{jws}{name={JWS}, sort={JWS},
  description={\emph{JSON Web Signature} : second format de token de Keystone, signé avec une clé asymétrique (ES256) ; son contenu n'est pas opaque}}

\newglossaryentry{ldap}{name={LDAP}, sort={LDAP},
  description={\emph{Lightweight Directory Access Protocol} : protocole d'accès à un annuaire (par exemple Active Directory) ; Keystone peut y lire utilisateurs et groupes, en lecture seule}}

\newglossaryentry{mapping}{name={mapping}, sort={mapping},
  description={Ensemble de règles qui traduisent les attributs fournis par un fournisseur d'identité en utilisateur, groupes et projets Keystone}}

\newglossaryentry{mfa}{name={authentification multifacteur (MFA)}, sort={authentification multifacteur}, text={authentification multifacteur},
  description={Authentification qui exige plusieurs méthodes successives, par exemple un mot de passe puis un code TOTP}}

\newglossaryentry{middleware}{name={middleware}, sort={middleware},
  description={Composant WSGI placé devant l'application d'un service ; \texttt{keystonemiddleware} (module \texttt{auth\_token}) y valide le token de chaque requête auprès de Keystone}}

\newglossaryentry{oidc}{name={OpenID Connect (OIDC)}, sort={OpenID Connect}, text={OpenID Connect},
  description={Protocole de fédération au format JSON, construit sur OAuth 2.0}}

\newglossaryentry{policy}{name={policy}, sort={policy}, plural={policies},
  description={Règle d'accès évaluée par chaque service (bibliothèque \texttt{oslo.policy}) pour décider si les rôles présents dans le token autorisent l'action demandée}}

\newglossaryentry{portee}{name={portée (scope)}, sort={portee}, text={portée}, plural={portées},
  description={Périmètre sur lequel un token confère des droits : un projet, un domaine ou le système entier ; un token non délimité n'a aucune portée}}

\newglossaryentry{projet}{name={projet (tenant)}, sort={projet}, text={projet}, plural={projets},
  description={Conteneur qui regroupe ou isole des ressources ; « tenant » est l'ancien nom du projet}}

\newglossaryentry{rbac}{name={RBAC}, sort={RBAC},
  description={\emph{Role-Based Access Control} : contrôle d'accès fondé sur les rôles attribués aux utilisateurs et aux groupes}}

\newglossaryentry{region}{name={région}, sort={region},
  description={Division générale d'un déploiement OpenStack, qui peut contenir des sous-régions ; chaque endpoint du catalogue est rattaché à une région}}

\newglossaryentry{role}{name={rôle}, sort={role},
  description={Ensemble de droits attribué à un utilisateur ou à un groupe sur un projet, un domaine ou le système ; les services l'interprètent pour décider des accès}}

\newglossaryentry{saml2}{name={SAML 2.0}, sort={SAML},
  description={\emph{Security Assertion Markup Language} : protocole de fédération au format XML, courant dans les universités et les entreprises}}

\newglossaryentry{sp}{name={fournisseur de services (SP)}, sort={fournisseur de services}, text={fournisseur de services}, plural={fournisseurs de services},
  description={Dans une fédération, service qui s'appuie sur un fournisseur d'identité ; ici Keystone, dont les tokens servent aux autres services OpenStack}}

\newglossaryentry{sso}{name={SSO}, sort={SSO},
  description={\emph{Single Sign-On}, authentification unique : une seule connexion donne accès à plusieurs services}}

\newglossaryentry{token}{name={token (jeton)}, sort={token}, text={token}, plural={tokens},
  description={Chaîne de caractères qui prouve l'identité d'un utilisateur et donne accès aux API ; valide pour une durée limitée, elle peut être révoquée à tout moment}}

\newglossaryentry{totp}{name={TOTP}, sort={TOTP},
  description={\emph{Time-based One-Time Password} : code à usage unique de six chiffres, calculé à partir d'un secret partagé et de l'heure ; second facteur pris en charge par Keystone}}

\newglossaryentry{trust}{name={trust}, sort={trust},
  description={Délégation par un utilisateur (le \emph{trustor}) d'une partie de ses rôles sur un projet à un autre utilisateur (le \emph{trustee})}}

\newglossaryentry{utilisateur}{name={utilisateur}, sort={utilisateur},
  description={Représentation numérique d'une personne, d'un système ou d'un service qui utilise les services OpenStack}}

\newglossaryentry{wsgi}{name={WSGI}, sort={WSGI},
  description={\emph{Web Server Gateway Interface} : interface standard entre un serveur web et une application Python ; Keystone s'exécute derrière Apache grâce à \texttt{mod\_wsgi}}}

% ---- Nova et Placement ------------------------------------------------
\newglossaryentry{instance}{name={instance}, sort={instance},
  description={Machine virtuelle créée et gérée par Nova ; on parle aussi de serveur (\texttt{openstack server})}}

\newglossaryentry{flavor}{name={gabarit (flavor)}, sort={gabarit}, text={gabarit},
  description={Modèle de taille d'une instance : nombre de vCPU, mémoire, disque racine, disque éphémère, swap, et propriétés supplémentaires (\emph{extra specs})}}

\newglossaryentry{noeudcalcul}{name={nœud de calcul (compute node)}, sort={noeud de calcul}, text={nœud de calcul},
  description={Serveur physique qui héberge les instances : il exécute \texttt{nova-compute} et un hyperviseur}}

\newglossaryentry{ordonnanceur}{name={ordonnanceur (scheduler)}, sort={ordonnanceur}, text={ordonnanceur},
  description={Composant de Nova (\texttt{nova-scheduler}) qui choisit le nœud de calcul où placer une instance}}

\newglossaryentry{conductor}{name={conductor}, sort={conductor},
  description={Composant de Nova (\texttt{nova-conductor}) qui coordonne les opérations longues (construction, redimensionnement, migration) et sert d'intermédiaire entre les nœuds de calcul et la base de données}}

\newglossaryentry{cellule}{name={cellule (cell)}, sort={cellule}, text={cellule},
  description={Groupe de nœuds de calcul de Nova qui possède sa propre base de données et sa propre file de messages ; tout déploiement a au moins une cellule}}

\newglossaryentry{libvirt}{name={libvirt}, sort={libvirt},
  description={Bibliothèque et service de pilotage des hyperviseurs (KVM/QEMU en particulier) que \texttt{nova-compute} utilise pour créer et contrôler les machines virtuelles}}

\newglossaryentry{busmessages}{name={bus de messages}, sort={bus de messages}, text={bus de messages},
  description={File de messages (RabbitMQ en général) par laquelle les composants d'un service échangent des appels distants (RPC, via \texttt{oslo.messaging})}}

\newglossaryentry{microversion}{name={microversion}, sort={microversion},
  description={Numéro de version fin d'une API OpenStack (pour Nova : 2.1, 2.2, 2.3 et ainsi de suite) : le client demande la version dont il a besoin, ce qui garde l'API compatible}}

\newglossaryentry{migrationchaud}{name={migration à chaud (live migration)}, sort={migration a chaud}, text={migration à chaud},
  description={Déplacement d'une instance en cours d'exécution d'un nœud de calcul vers un autre, sans l'arrêter}}

\newglossaryentry{surallocation}{name={surallocation (overcommit)}, sort={surallocation}, text={surallocation},
  description={Attribution de plus de ressources virtuelles qu'il n'y en a de physiques ; elle se règle par un ratio d'allocation}}

\newglossaryentry{agregat}{name={agrégat d'hôtes (host aggregate)}, sort={agregat}, text={agrégat d'hôtes},
  description={Groupe de nœuds de calcul défini par l'administrateur, avec des métadonnées ; il sert à l'ordonnancement et à définir des zones de disponibilité}}

\newglossaryentry{zonedispo}{name={zone de disponibilité}, sort={zone de disponibilite}, text={zone de disponibilité},
  description={Groupe de nœuds de calcul que l'utilisateur peut viser à la création d'une instance (par exemple une salle ou une baie) ; c'est un agrégat d'hôtes visible des utilisateurs}}

\newglossaryentry{fournisseur}{name={fournisseur de ressources (resource provider)}, sort={fournisseur de ressources}, text={fournisseur de ressources},
  description={Entité de Placement qui offre un inventaire de ressources consommables, par exemple un nœud de calcul ou un pool de stockage partagé}}

\newglossaryentry{inventaire}{name={inventaire}, sort={inventaire},
  description={Quantité d'une classe de ressources offerte par un fournisseur : total, réservé, ratio d'allocation, taille minimale, maximale et pas d'allocation}}

\newglossaryentry{classeressources}{name={classe de ressources}, sort={classe de ressources}, text={classe de ressources},
  description={Type de ressource quantitative suivie par Placement : \texttt{VCPU}, \texttt{MEMORY\_MB}, \texttt{DISK\_GB}, ou une classe personnalisée préfixée par \texttt{CUSTOM\_}}}

\newglossaryentry{allocation}{name={allocation}, sort={allocation},
  description={Quantité de ressources d'un fournisseur attribuée à un consommateur ; elle diminue d'autant la capacité disponible}}

\newglossaryentry{consommateur}{name={consommateur (consumer)}, sort={consommateur}, text={consommateur},
  description={Entité qui détient des allocations dans Placement, identifiée par un UUID : pour Nova, une instance ou une migration}}

\newglossaryentry{trait}{name={trait}, sort={trait},
  description={Caractéristique qualitative d'un fournisseur (par exemple \texttt{HW\_CPU\_X86\_AVX2} ou \texttt{STORAGE\_DISK\_SSD}) ; les traits non standard commencent par \texttt{CUSTOM\_}}}

\newglossaryentry{claim}{name={réservation (claim)}, sort={reservation}, text={réservation},
  description={Création d'une allocation dans Placement pour une instance, faite par l'ordonnanceur une fois l'hôte choisi ; si elle échoue, c'est qu'un autre consommateur a pris les ressources entre-temps}}

% ---- Glance ----------------------------------------------------------
\newglossaryentry{image}{name={image}, sort={image},
  description={Fichier disque d'une machine virtuelle (système installé) qu'on duplique pour créer des instances ; Glance en conserve les données et les métadonnées}}

\newglossaryentry{qcow2}{name={qcow2}, sort={qcow2},
  description={Format de disque de QEMU à copie sur écriture : le fichier ne grossit qu'avec les données écrites et peut s'appuyer sur un fichier de base}}

\newglossaryentry{store}{name={store (backend de stockage)}, sort={store}, text={store},
  description={Système où Glance garde les données des images : système de fichiers, Ceph (rbd), Swift, Cinder, VMware ; en lecture seule pour http}}

% ---- Neutron ---------------------------------------------------------
\newglossaryentry{reseau}{name={réseau (network)}, sort={reseau}, text={réseau}, plural={réseaux},
  description={Dans Neutron, domaine de diffusion de niveau 2 isolé, auquel on rattache des ports}}

\newglossaryentry{sousreseau}{name={sous-réseau (subnet)}, sort={sous-reseau}, text={sous-réseau}, plural={sous-réseaux},
  description={Plage d'adresses IP associée à un réseau Neutron, avec sa passerelle et son serveur DHCP éventuel}}

\newglossaryentry{port}{name={port}, sort={port},
  description={Point de connexion d'un équipement (en général la carte réseau d'une instance) à un réseau Neutron ; il porte une adresse MAC, des adresses IP et des groupes de sécurité}}

\newglossaryentry{routeur}{name={routeur (router)}, sort={routeur}, text={routeur},
  description={Objet Neutron qui relie des réseaux entre eux et vers l'extérieur, avec traduction d'adresses (NAT)}}

\newglossaryentry{groupesecurite}{name={groupe de sécurité}, sort={groupe de securite}, text={groupe de sécurité}, plural={groupes de sécurité},
  description={Ensemble de règles de pare-feu appliqué aux ports : par défaut tout le trafic entrant est refusé et tout le trafic sortant autorisé ; les règles sont avec état}}

\newglossaryentry{ipflottante}{name={IP flottante (floating IP)}, sort={IP flottante}, text={IP flottante},
  description={Adresse publique associée à la demande à un port d'instance, et traduite en son adresse privée par le routeur}}

\newglossaryentry{provider}{name={réseau provider}, sort={reseau provider}, text={réseau provider},
  description={Réseau créé par l'administrateur et rattaché directement à un réseau physique existant (flat ou VLAN)}}

\newglossaryentry{selfservice}{name={réseau self-service}, sort={reseau self-service}, text={réseau self-service},
  description={Réseau virtuel (overlay) créé par un projet sans droits d'administrateur ; il a besoin d'un routeur pour sortir}}

\newglossaryentry{overlay}{name={overlay (réseau superposé)}, sort={overlay}, text={overlay},
  description={Réseau virtuel dont les trames sont encapsulées (VXLAN, Geneve) dans des paquets IP du réseau physique}}

\newglossaryentry{vlan}{name={VLAN}, sort={VLAN},
  description={Réseau local virtuel 802.1Q identifié par un numéro de 1 à 4094 ; la limite de 4094 réseaux restreint son usage pour les projets}}

\newglossaryentry{vxlan}{name={VXLAN, Geneve}, sort={VXLAN}, text={VXLAN},
  description={Protocoles de tunnel qui encapsulent des trames Ethernet dans des paquets UDP/IP ; ils servent aux réseaux overlay et n'ont pas la limite des VLAN}}

\newglossaryentry{ml2}{name={ML2 (Modular Layer 2)}, sort={ML2}, text={ML2},
  description={Plugin de base de Neutron : il orchestre des \emph{type drivers} (flat, VLAN, VXLAN\dots) et des \emph{mechanism drivers} (Open vSwitch, OVN\dots)}}

\newglossaryentry{ovs}{name={Open vSwitch (OVS)}, sort={Open vSwitch}, text={Open vSwitch},
  description={Commutateur virtuel programmable : il relie sur chaque nœud les interfaces des instances, les tunnels et les interfaces physiques, selon des règles (\emph{flows}) OpenFlow}}

\newglossaryentry{ovn}{name={OVN (Open Virtual Network)}, sort={OVN}, text={OVN},
  description={Système de réseau virtuel bâti sur Open vSwitch, avec deux bases de données et un démon par nœud ; il distribue le routage, le DHCP et les règles de sécurité}}

\newglossaryentry{agentneutron}{name={agent Neutron}, sort={agent Neutron}, text={agent},
  description={Processus qui, sur un nœud, traduit les objets Neutron en configuration réelle : agent L2 (Open vSwitch), L3 (routeurs), DHCP, métadonnées}}

\newglossaryentry{nat}{name={NAT}, sort={NAT},
  description={Traduction d'adresses réseau : le routeur remplace l'adresse source ou destination d'un paquet (SNAT pour la sortie, DNAT pour l'entrée)}}

% ---- Cinder -----------------------------------------------------------
\newglossaryentry{volume}{name={volume}, sort={volume},
  description={Disque virtuel persistant fourni par Cinder : il a une taille et un type, et peut être attaché à une instance, détaché, puis rattaché ailleurs sans perdre ses données}}

\newglossaryentry{snapshot}{name={snapshot}, sort={snapshot},
  description={Image ponctuelle d'un volume, conservée sur le même backend ; elle sert au retour arrière et au clonage, mais disparaît avec le backend}}

\newglossaryentry{sauvegarde}{name={sauvegarde (backup)}, sort={sauvegarde}, text={sauvegarde},
  description={Copie complète ou incrémentale d'un volume vers un dépôt distinct (Swift par défaut) ; elle survit à la perte du backend}}

\newglossaryentry{typevolume}{name={type de volume}, sort={type de volume}, plural={types de volume},
  description={Étiquette qui porte des propriétés (\emph{extra specs}) : backend visé, qualité de service, accès multiple ; Cinder s'en sert pour choisir le backend d'un volume}}

\newglossaryentry{backendcinder}{name={backend (Cinder)}, sort={backend}, text={backend}, plural={backends},
  description={Système de stockage qui héberge réellement les volumes (LVM, Ceph, NFS, baie SAN) ; \texttt{cinder-volume} le pilote par un driver}}

\newglossaryentry{qos}{name={QoS (qualité de service)}, sort={QoS}, text={QoS},
  description={Limite de débit appliquée à un volume, en opérations ou en octets par seconde ; elle s'associe à un type de volume}}

\newglossaryentry{osbrick}{name={os-brick}, sort={os-brick},
  description={Bibliothèque qui connecte un hôte à un volume (iSCSI, RBD, NFS\dots) ; Nova l'utilise sur le nœud de calcul pour attacher les volumes}}

% ---- Horizon ----------------------------------------------------------
\newglossaryentry{django}{name={Django}, sort={Django},
  description={Cadriciel Web en Python sur lequel est écrit Horizon ; il fournit les sessions, l'authentification par des backends et le système de gabarits}}

\newglossaryentry{dashboard}{name={dashboard (tableau de bord)}, sort={dashboard}, text={dashboard},
  description={Entrée de la navigation principale d'Horizon (\emph{Project}, \emph{Admin}, \emph{Identity}, \emph{Settings}) ; il regroupe des groupes de panneaux}}

\newglossaryentry{panneau}{name={panneau (panel)}, sort={panneau}, text={panneau}, plural={panneaux},
  description={Page de la navigation latérale d'un dashboard, déclarée dans un fichier \texttt{panel.py} ; elle a ses routes, ses vues et ses gabarits}}

\newglossaryentry{memcached}{name={memcached}, sort={memcached},
  description={Service de cache en mémoire, partagé entre processus et serveurs ; Horizon l'utilise couramment pour conserver les sessions}}

% ---- Octavia ----------------------------------------------------------
\newglossaryentry{loadbalancer}{name={répartiteur de charge (load balancer)}, sort={repartiteur de charge}, text={répartiteur de charge}, plural={répartiteurs de charge},
  description={Service qui reçoit le trafic sur une adresse unique (la VIP) et le distribue entre plusieurs serveurs ; c'est l'objet principal que l'utilisateur crée dans Octavia}}

\newglossaryentry{listener}{name={listener}, sort={listener}, plural={listeners},
  description={Point d'écoute d'un répartiteur : un protocole et un port (par exemple HTTP, 80) sur la VIP ; il envoie le trafic reçu vers un pool}}

\newglossaryentry{poolocta}{name={pool (Octavia)}, sort={pool}, text={pool}, plural={pools},
  description={Groupe de serveurs qui reçoivent le trafic d'un listener, avec un algorithme de répartition et, au besoin, un health monitor}}

\newglossaryentry{memberocta}{name={member (Octavia)}, sort={member}, text={member}, plural={members},
  description={Serveur d'un pool, repéré par une adresse IP et un port ; il peut être une instance Nova ou n'importe quelle adresse joignable}}

\newglossaryentry{healthmonitor}{name={health monitor}, sort={health monitor},
  description={Sonde que le répartiteur envoie à chaque member (TCP, HTTP, HTTPS, PING\dots) ; un member qui échoue est retiré de la distribution jusqu'à ce qu'il réponde de nouveau}}

\newglossaryentry{amphore}{name={amphore}, sort={amphore}, plural={amphores},
  description={Instance Nova, créée par Octavia, qui exécute HAProxy et porte la VIP ; c'est elle qui écoute et répartit réellement le trafic}}

\newglossaryentry{vip}{name={VIP (adresse virtuelle)}, sort={VIP}, text={VIP},
  description={Adresse IP du répartiteur de charge, celle que les clients utilisent ; elle est portée par un port Neutron et peut recevoir une adresse flottante}}

\newglossaryentry{l7}{name={politique L7 (L7 policy)}, sort={politique L7}, text={politique L7}, plural={politiques L7},
  description={Règle de routage au niveau applicatif (chemin d'URL, en-tête, nom d'hôte\dots) qui redirige une requête vers un autre pool ou vers une URL}}

\newglossaryentry{haproxy}{name={HAProxy}, sort={HAProxy},
  description={Logiciel libre de répartition de charge TCP et HTTP ; l'amphore de référence d'Octavia l'exécute}}

% ---- Swift ------------------------------------------------------------
\newglossaryentry{objetswift}{name={objet (Swift)}, sort={objet}, text={objet}, plural={objets},
  description={Donnée stockée dans Swift : le contenu d'un fichier et ses métadonnées, adressé par son chemin \texttt{/compte/conteneur/objet}}}

\newglossaryentry{conteneur}{name={conteneur (Swift)}, sort={conteneur}, text={conteneur}, plural={conteneurs},
  description={Espace de noms d'un compte Swift qui regroupe des objets ; il porte les droits d'accès, la politique de stockage et d'autres métadonnées}}

\newglossaryentry{proxyswift}{name={proxy (Swift)}, sort={proxy}, text={proxy}, plural={proxys},
  description={Service \texttt{swift-proxy-server} : seul point d'entrée de l'API, il authentifie, consulte le ring et relaie les requêtes vers les serveurs de stockage}}

\newglossaryentry{ring}{name={ring}, sort={ring}, plural={rings},
  description={Fichier qui associe chaque partition à des disques précis ; Swift en tient un pour les comptes, un pour les conteneurs et un par politique de stockage d'objets}}

\newglossaryentry{partition}{name={partition (Swift)}, sort={partition}, text={partition}, plural={partitions},
  description={Unité de répartition du ring : un intervalle des empreintes MD5 des chemins ; le ring place ses répliques sur des disques de zones différentes}}

\newglossaryentry{handoff}{name={handoff}, sort={handoff}, plural={handoffs},
  description={Disque de remplacement choisi par le ring lorsqu'un disque normal est en panne ; la réplication ramène ensuite les données à leur place}}

\newglossaryentry{erasurecoding}{name={codage d'effacement (erasure coding)}, sort={codage d'effacement}, text={codage d'effacement},
  description={Méthode de redondance qui découpe un objet en fragments de données et de parité ; Swift l'offre par une politique de stockage, pour un surcoût moindre que la triple réplication}}

\newglossaryentry{politiquestockage}{name={politique de stockage (storage policy)}, sort={politique de stockage}, text={politique de stockage}, plural={politiques de stockage},
  description={Choix, par conteneur, d'un ring d'objets et donc du nombre de répliques, du type de disques ou du codage d'effacement}}

% ---- Heat -------------------------------------------------------------
\newglossaryentry{modeleheat}{name={modèle (template)}, sort={modele}, text={modèle}, plural={modèles},
  description={Fichier qui décrit l'infrastructure voulue (ressources, paramètres, sorties) ; Heat le lit pour créer une stack. Son format natif est HOT, en YAML}}

\newglossaryentry{stack}{name={stack}, sort={stack}, plural={stacks},
  description={Ensemble de ressources créé par Heat à partir d'un modèle ; c'est l'unité qu'on crée, met à jour et supprime}}

\newglossaryentry{ressourceheat}{name={ressource (Heat)}, sort={ressource}, text={ressource}, plural={ressources},
  description={Élément d'infrastructure d'une stack (réseau, serveur, volume\dots), instancié d'après un type de la forme \texttt{OS::Service::Type}}}

\newglossaryentry{hot}{name={HOT (Heat Orchestration Template)}, sort={HOT}, text={HOT},
  description={Format de modèle natif de Heat, écrit en YAML ; sa version (\texttt{heat\_template\_version}) détermine les fonctions disponibles}}

\newglossaryentry{environnementheat}{name={environnement (Heat)}, sort={environnement}, text={environnement},
  description={Fichier passé avec \texttt{-e} qui fournit des valeurs de paramètres et peut remplacer des types de ressources (\texttt{resource\_registry})}}

\newglossaryentry{graphedeps}{name={graphe de dépendances}, sort={graphe de dependances}, text={graphe de dépendances},
  description={Graphe orienté que Heat construit à partir des références entre ressources ; il fixe l'ordre de création, et la suppression suit l'ordre inverse}}

% ---- Ironic -----------------------------------------------------------
\newglossaryentry{baremetal}{name={bare metal}, sort={bare metal}, text={bare metal},
  description={Serveur physique utilisé directement par un seul locataire, sans hyperviseur ; Ironic le provisionne et Nova le présente comme une instance}}

\newglossaryentry{noeudironic}{name={nœud (Ironic)}, sort={noeud (Ironic)}, text={nœud}, plural={nœuds},
  description={Enregistrement d'Ironic qui représente un serveur physique : son type de matériel, ses identifiants de gestion, ses ports, sa classe de ressources et son état de provisionnement}}

\newglossaryentry{portironic}{name={port (Ironic)}, sort={port (Ironic)}, text={port}, plural={ports},
  description={Carte réseau d'un nœud, déclarée dans Ironic par son adresse MAC ; à ne pas confondre avec le port Neutron, qui lui est lié lors du raccordement au réseau}}

\newglossaryentry{conductorironic}{name={conductor (Ironic)}, sort={conductor (Ironic)}, text={conductor}, plural={conductors},
  description={Composant d'Ironic (\texttt{ironic-conductor}) qui agit sur les serveurs physiques : alimentation, déploiement, nettoyage, inspection ; chaque conductor gère une partie des nœuds}}

\newglossaryentry{bmc}{name={BMC (contrôleur de gestion)}, sort={BMC}, text={BMC},
  description={\emph{Baseboard Management Controller} : contrôleur intégré au serveur, doté de sa propre adresse réseau ; il reçoit les ordres d'alimentation et de démarrage (IPMI, Redfish)}}

\newglossaryentry{ipmi}{name={IPMI}, sort={IPMI}, text={IPMI},
  description={\emph{Intelligent Platform Management Interface} : protocole ancien de gestion à distance par le BMC (alimentation, périphérique de démarrage) ; Ironic l'utilise avec la commande \texttt{ipmitool}}}

\newglossaryentry{redfish}{name={Redfish}, sort={Redfish}, text={Redfish},
  description={Protocole de gestion à distance des serveurs par API REST, proposé par le BMC ; Ironic l'utilise avec le type de matériel \texttt{redfish}}}

\newglossaryentry{pxe}{name={PXE (démarrage par le réseau)}, sort={PXE}, text={PXE},
  description={\emph{Preboot Execution Environment} : démarrage par le réseau, où un serveur DHCP indique au serveur le programme à charger et un serveur TFTP (ou HTTP) le fournit ; iPXE en est une variante}}

\newglossaryentry{ipa}{name={IPA (Ironic Python Agent)}, sort={IPA}, text={IPA},
  description={Agent Python qui s'exécute dans un ramdisk sur le serveur : il donne au conductor un accès à distance et réalise de l'intérieur l'écriture de l'image, le nettoyage et l'inspection}}

\newglossaryentry{hardwaretype}{name={type de matériel (hardware type)}, sort={type de materiel}, text={type de matériel}, plural={types de matériel},
  description={Ensemble d'interfaces matérielles qui définit comment Ironic pilote une famille de serveurs (\texttt{ipmi}, \texttt{redfish} ou un type propre à un constructeur) ; on parle aussi de pilote (\emph{driver})}}

\newglossaryentry{interfacemat}{name={interface matérielle (hardware interface)}, sort={interface materielle}, text={interface matérielle}, plural={interfaces matérielles},
  description={Brique d'un type de matériel qui gère un aspect du serveur : alimentation, démarrage, déploiement, réseau, inspection, RAID\dots}}

\newglossaryentry{nettoyage}{name={nettoyage (cleaning)}, sort={nettoyage}, text={nettoyage},
  description={Étapes qu'Ironic exécute sur un nœud entre deux locataires, ou avant le premier : effacement des disques, remise du matériel dans une configuration connue ; l'état \texttt{cleaning} les regroupe}}

\newglossaryentry{ramdisk}{name={ramdisk de déploiement}, sort={ramdisk}, text={ramdisk},
  description={Petit système Linux chargé en mémoire par démarrage réseau, qui exécute l'agent IPA pour déployer, nettoyer ou inspecter un nœud}}

\newglossaryentry{etatprov}{name={état de provisionnement (provision state)}, sort={etat de provisionnement}, text={état de provisionnement}, plural={états de provisionnement},
  description={Étape du cycle de vie d'un nœud dans la machine à états d'Ironic (\texttt{enroll}, \texttt{manageable}, \texttt{available}, \texttt{active}\dots)}}
